
\documentclass[11pt]{book}
\usepackage[a4paper,margin=1.15in]{geometry}
\usepackage[T1]{fontenc}
\usepackage[utf8]{inputenc}
\usepackage{lmodern}
\usepackage{microtype}
\usepackage{amsmath,amssymb,mathtools}
\usepackage{longtable,array,booktabs}
\usepackage{enumitem}
\usepackage{hyperref}
\hypersetup{
  hypertexnames=false,
  hidelinks,
  pdfauthor={Luigi Bianchi},
  pdftitle={Lectures on the Theory of Finite Continuous Groups of Transformations -- English -- manuscript pages 1--116},
  pdfsubject={Cumulative English translation from the published handwritten 1903 volume}
}
\setlength{\parindent}{1.3em}
\setlength{\parskip}{0.2em}
\setlength{\emergencystretch}{4em}
\hbadness=10000
\newcommand{\sen}{\operatorname{sen}}
\newcommand{\dif}{\,d}
\newcommand{\entry}[3]{#1 & #2 & #3\\}
\begin{document}

\begin{titlepage}
\centering
{\Large LUIGI BIANCHI\par}
\vspace{0.4em}
{\small PROFESSOR OF THE ROYAL UNIVERSITY OF PISA\par}
\vfill
{\Huge LECTURES\par}
\vspace{1.5em}
{\large ON THE\par}
\vspace{1.5em}
{\Large Theory of finite continuous groups of transformations\par}
\vspace{2em}
{\large Year 1902--1903\par}
\vfill
{\Large PISA\par}
{\large ENRICO SPOERRI\par}
{\small BOOKSELLER-PUBLISHER\par}
\vspace{0.5em}
{\large 1903\par}
\end{titlepage}

\frontmatter
\chapter*{Table of chapters}
\footnotesize
\begin{longtable}{@{}>{\raggedright\arraybackslash}p{0.42in} >{\raggedright\arraybackslash}p{4.60in} >{\raggedleft\arraybackslash}p{0.32in}@{}}
\toprule
\textbf{§} & \textbf{Title} & \textbf{p.}\\
\midrule
\endfirsthead
\toprule
\textbf{§} & \textbf{Title} & \textbf{p.}\\
\midrule
\endhead
\multicolumn{3}{c}{\textsc{Chapter I}}\\
\entry{}{\textit{Generalities on finite continuous groups of transformations. -- Fundamental differential equations. -- Infinitesimal transformations as generators of the group.}}{}
\entry{1}{Transformations and their composition}{1}
\entry{2}{Series $\infty^r$ of transformations}{5}
\entry{3}{Criterion for recognizing the number of essential parameters}{10}
\entry{4}{Definition of finite continuous groups}{16}
\entry{5}{Examples of continuous groups}{20}
\entry{6}{Solvability of the equations $c_i=\varphi_i(a,b)$ with respect to the $a$ or to the $b$}{24}
\entry{7}{Fundamental differential equations}{27}
\entry{8}{Some properties of the functions $\psi_{sh}(a)$, $\xi_{si}(x')$}{31}
\entry{9}{One-parameter groups $G_1$}{34}
\entry{10}{Infinitesimal transformation of a $G_1$}{38}
\entry{11}{Change of variables in the infinitesimal transformations}{42}
\entry{12}{Series $\infty^r$ of transformations and their infinitesimal transformations}{43}
\entry{13}{Particular case of a group}{47}
\entry{14}{New derivation of the fundamental equations (A)}{50}
\entry{15}{The $r$ infinitesimal transformations $X_\lambda f=\sum\limits_{i=1}^n \xi_{\lambda i}(x)\dfrac{\partial f}{\partial x_i}$}{52}
\entry{16}{Fundamental theorem for the series of transformations satisfying (A)}{55}
\addlinespace[0.5em]
\multicolumn{3}{c}{\textsc{Chapter II}}\\
\entry{}{\textit{The three fundamental theorems of the theory of groups. -- Mixed groups.}}{}
\entry{17}{The first fundamental theorem}{58}
\entry{18}{Further cases}{61}
\entry{19}{Complete system of partial differential equations equivalent to the system (A)}{63}
\entry{20--21}{Integrability conditions $(\alpha)$ and $(\beta)$}{65}
\entry{22}{Schur's proof of the first fundamental theorem}{69}
\entry{23}{The second fundamental theorem (principal theorem)}{72}
\entry{24}{The relations among the composition constants $C_{iks}$}{77}
\entry{25}{Property of the symbol $|UV|:=\sum\limits_{i,k,s=1}^{r}C_{iks}y_s\dfrac{\partial U}{\partial y_i}\dfrac{\partial V}{\partial y_k}$}{80}
\entry{26}{Change of variables}{83}
\entry{27}{The third fundamental theorem}{86}
\entry{28}{Maurer's equations for the $\varphi_{ik}$}{89}
\entry{29}{Schur's proof of the second and third fundamental theorems}{92}
\entry{30}{Composition of two series $\infty^r$ of transformations}{95}
\entry{31}{Series $\infty^r$ of transformations in which the products of $m$ of their transformations belong to the series}{100}
\entry{32}{Mixed (or complex) groups}{102}
\addlinespace[0.5em]
\multicolumn{3}{c}{\textsc{Chapter III}}\\
\entry{}{\textit{Equations defining a group. -- Order of the infinitesimal transformations.}}{}
\entry{33}{Elimination of arbitrary constants from a system of functions}{107}
\entry{34}{Mayer systems of partial differential equations and their properties}{111}
\entry{35}{The equations defining a group $G_r$}{116}
\entry{36}{Another way to define a continuous group. -- Note on infinite groups}{119}
\entry{37}{Order of an infinitesimal transformation}{124}
\entry{38}{Distribution of the orders of the infinitesimal transformations in a group}{127}
\entry{39}{Calculation of the numbers $r$ from the equations defining the group}{132}
\entry{40}{Subgroups generated by the transformations of higher orders}{137}
\entry{41}{Example of the group of motions}{144}
\addlinespace[0.5em]
\multicolumn{3}{c}{\textsc{Chapter IV}}\\
\entry{}{\textit{Functions and varieties invariant with respect to a group. -- Transitive and intransitive groups. -- Imprimitive groups. -- Systatic and asystatic groups.}}{}
\entry{42}{Functions invariant with respect to a group}{148}
\entry{43}{Digression on varieties in a space $S_n$ of $n$ dimensions}{150}
\entry{44}{Varieties invariant with respect to a group $G_1$}{154}
\entry{45}{Varieties invariant with respect to a group $G_r$}{159}
\entry{46}{Dimension of the minimum invariant varieties}{161}
\entry{47}{Determination of all minimum invariant varieties}{164}
\entry{48}{Transitive groups}{169}
\entry{49}{Invariants of intransitive groups}{171}
\entry{50}{Various examples}{174}
\entry{51}{Imprimitive groups}{178}
\entry{52}{Complete systems admitting a group}{182}
\entry{53}{Systatic and asystatic groups}{186}
\entry{54}{Imprimitivity of a systatic group}{190}
\entry{55}{Criterion for systaticity deduced from the equations of definition}{193}
\addlinespace[0.5em]
\multicolumn{3}{c}{\textsc{Chapter V}}\\
\entry{}{\textit{Invariant series of infinitesimal transformations. -- Invariant subgroups and composition series of a group.}}{}
\entry{56}{Invariant series of infinitesimal transformations}{198}
\entry{57}{Invariant series with respect to a group}{201}
\entry{58}{Finite transformations of the series $\sum\limits_h^{1,\ldots,r} e_h X_h f$}{207}
\entry{59}{Permutable transformations}{209}
\entry{60}{Invariant subgroups. -- Simple groups}{212}
\entry{61}{Composition series. -- Integrable groups}{215}
\entry{62}{Derived group}{218}
\entry{63}{Integrability conditions of a group}{221}
\entry{64}{Product of two permutable groups}{224}
\entry{65}{Theorem concerning two maximum invariant subgroups}{226}
\entry{66}{Jordan-Lie theorem on the indices of the composition series}{229}
\addlinespace[0.5em]
\multicolumn{3}{c}{\textsc{Chapter VI}}\\
\entry{}{\textit{The adjoint group and the subgroups of a group. -- Two-parameter subgroups. -- Characteristic equation. -- Composition or structure of a group. -- Genus of a group.}}{}
\entry{67}{Permutations among the transformations of a group induced by the transformations themselves}{232}
\entry{68}{Canonical form of the equations of a group}{235}
\entry{69}{The adjoint group and its infinitesimal transformations}{239}
\entry{70}{Number of parameters of the adjoint group}{243}
\entry{71}{Search for the subgroups of a group}{247}
\entry{72}{Affine subgroups}{249}
\entry{73}{Representation in space of $r$ dimensions}{252}
\entry{74}{Example of the group of motions in the plane}{256}
\entry{75}{Application to the projective group on one variable}{262}
\entry{76}{Two-parameter subgroups of a group $G_r$}{266}
\entry{77}{The characteristic equation}{271}
\entry{78}{Properties of the characteristic equation}{276}
\entry{79}{The coefficients of $\Delta(\rho)$ as invariants of the adjoint group}{279}
\entry{80}{Composition (or structure) of a group}{282}
\entry{81}{The relations among the $C_{iks}$, $C'_{iks}$ as determining a group}{285}
\entry{82}{Genus (rank) of a group}{287}
\addlinespace[0.5em]
\multicolumn{3}{c}{\textsc{Chapter VII}}\\
\entry{}{\textit{First investigations on similarity of groups. -- Transformations permutable with all the transformations of a group. -- Reciprocal simply transitive groups.}}{}
\entry{83}{Similarity of groups}{291}
\entry{84}{Case in which the condition of equal composition is sufficient}{296}
\entry{85}{Case of Abelian groups and simply transitive groups}{302}
\entry{86}{General conditions necessary for similarity}{303}
\entry{87}{Case of asystatic groups}{306}
\entry{88}{Similarity of a group with itself}{309}
\entry{89}{Transformations permutable with the individual transformations of a group}{312}
\entry{90}{Reciprocal simply transitive groups}{317}
\entry{91}{Properties of pairs of reciprocal groups}{320}
\entry{92}{Direct construction of the reciprocal group $\Gamma_n$}{325}
\entry{93}{Imprimitivity of a simply transitive group}{328}
\addlinespace[0.5em]
\multicolumn{3}{c}{\textsc{Chapter VIII}}\\
\entry{}{\textit{Isomorphism of groups referred to their infinitesimal transformations. -- Parametric groups. -- Isomorphism of groups in relation to finite transformations. -- Solution of the general problem of recognizing the similarity of two groups.}}{}
\entry{94}{Isomorphism of two groups referred to their infinitesimal transformations}{331}
\entry{95}{Properties of isomorphism}{337}
\entry{96}{Complementary group with respect to an invariant subgroup}{340}
\entry{97}{Composition series of two isomorphic groups. -- Hölder-Engel theorem}{345}
\entry{98}{Construction of groups isomorphic to primitive groups}{350}
\entry{99}{Definition of parametric groups}{354}
\entry{100}{Infinitesimal transformations of the (first) parametric group}{359}
\entry{101}{Equivalent parametric groups}{363}
\entry{102}{Holoedric isomorphism in relation to finite transformations}{366}
\entry{103}{Meriedric isomorphism}{372}
\entry{104}{Return to the conditions of similarity for two groups}{376}
\entry{105}{Similarity among transitive groups}{379}
\entry{106}{Similarity among intransitive groups}{384}
\entry{107}{New statement for the conditions of similarity}{391}
\entry{108}{Transformations permutable with a group}{394}
\entry{109}{Theorems on asystatic and systatic groups}{397}
\addlinespace[0.5em]
\multicolumn{3}{c}{\textsc{Chapter IX}}\\
\entry{}{\textit{Extended and prolonged groups. -- Differential invariants.}}{}
\entry{110}{Multiple transitivity}{401}
\entry{111}{Invariants of groups of points}{404}
\entry{112}{Groups prolonged to linear elements}{408}
\entry{113}{Infinitesimal transformations of the prolonged group}{412}
\entry{114}{Group induced by the stabilizer subgroup on the linear elements through a point}{415}
\entry{115}{Definition of differential invariants}{419}
\entry{116}{Prolonged transformations}{422}
\entry{117}{Prolonged groups}{426}
\entry{118}{Existence and calculation of differential invariants}{431}
\entry{119}{Differential invariants of groups on two variables}{433}
\entry{120}{Examples of calculation of differential invariants}{439}
\entry{121}{Differential invariants of surfaces with respect to motions}{443}
\addlinespace[0.5em]
\multicolumn{3}{c}{\textsc{Chapter X}}\\
\entry{}{\textit{Finite continuous groups on one and on two variables.}}{}
\entry{122}{Number of parameters of a group on one variable}{446}
\entry{123}{Reduction of groups on one variable to projective groups}{449}
\entry{124}{Homogeneous linear group $G_4$ on two variables}{453}
\entry{125}{Invariant subgroups and subgroups $G_1$ of $G_4$}{456}
\entry{126}{Two-parameter subgroups $G_2$ in $G_4$}{460}
\entry{127}{Subgroups $G_3$ of $G_4$ and recapitulation}{464}
\entry{128}{Criterion for distinguishing primitive groups of the plane from imprimitive ones}{466}
\entry{129}{Maximum order of the infinitesimal transformations in a primitive group}{471}
\entry{130}{Number of parameters of a primitive group in the plane}{474}
\entry{131}{Reduction of primitive groups to normal form in case A)}{478}
\entry{132}{Reduction to normal form in cases B) and C)}{482}
\entry{133}{The groups $G_8$, $G_6$, $G_5$ as projective groups}{486}
\entry{134}{Intransitive groups of the plane. -- Cases in which the group $G_r$ has one or two parameters}{490}
\entry{135}{Case in which the reduced group $G'_r$ has three parameters}{495}
\entry{136}{Imprimitive transitive groups. -- Case in which the shortened group $\Gamma'$ has one parameter}{500}
\entry{137}{Case in which $\Gamma'$ has two parameters}{505}
\entry{138}{Case in which $\Gamma'$ has three parameters}{511}
\addlinespace[0.5em]
\multicolumn{3}{c}{\textsc{Chapter XI}}\\
\entry{}{\textit{First applications of the theory of finite continuous groups to problems of integration of differential equations.}}{}
\entry{139}{Preliminaries}{518}
\entry{140}{First-order differential equations admitting a group $Xf$}{520}
\entry{141}{Relations between infinitesimal transformations and integrating factors}{524}
\entry{142}{Determination of the infinitesimal transformations admitted by a first-order differential equation}{527}
\entry{143}{Geometric interpretation of the integrating factors}{532}
\entry{144}{Case of first-order equations not solved with respect to the derivative}{536}
\entry{145}{Determination of the first-order equations that admit a given $G_1$}{538}
\entry{146}{Systems of differential equations and their infinitesimal transformations}{542}
\entry{147}{Properties of the infinitesimal transformations admitted by a differential system}{544}
\entry{148}{Integration of an $Af=0$ with known infinitesimal transformations}{549}
\entry{149}{Equations $Af=0$ with an integrable group $G_n$}{553}
\entry{150}{Reduction of an equation $Af=0$ with composite group to successive equations with simple groups}{556}
\entry{151}{Jacobi multipliers for an equation $Af=0$}{563}
\entry{152}{Relations between the multipliers and the infinitesimal transformations}{568}
\entry{153}{Integration of equations $Af=0$ in three variables with a known infinitesimal transformation}{572}
\entry{154}{Case in which two infinitesimal transformations of $Af=0$ are known}{574}
\entry{155}{Criterion for recognizing whether a second-order equation $y''=\omega(x,y,y')$ admits a group $Xf$}{580}
\entry{156}{Maximum number of infinitesimal transformations admitted by a second-order equation}{584}
\entry{157}{Second-order differential equations with a group $G_3$}{589}
\entry{158}{Second-order or higher-order differential equations with a given $G_1$}{592}
\entry{159}{Integration of a second-order differential equation with a known $G_1$}{595}
\entry{160}{Reduction to normal form of Abelian groups $G_2$ on two variables}{598}
\entry{161}{Case of a non-Abelian $G_2$}{602}
\entry{162}{Integration of a second-order differential equation with a known $G_2$}{605}
\addlinespace[0.5em]
\multicolumn{3}{c}{\textsc{Chapter XII}}\\
\entry{}{\textit{Applications to the theory of linear differential equations. -- Extension of the theories of Galois according to Picard and Vessiot.}}{}
\entry{163}{Criterion for linear dependence of several functions}{609}
\entry{164}{Fundamental systems of integrals of a linear differential equation}{614}
\entry{165}{Rational differential invariants of the total homogeneous linear group $L$}{617}
\entry{166}{Invariants of the subgroups of a group $G_r$}{623}
\entry{167}{Resolvents of linear differential equations}{628}
\entry{168}{Picard's resolvent and its fundamental properties}{634}
\entry{169}{Field of rationality and algebraic differential equations}{642}
\entry{170}{Group of rationality of a linear differential equation}{648}
\entry{171}{Fundamental property of the group of rationality}{654}
\entry{172}{Another definition of the group of rationality. -- Group of the general equation}{658}
\entry{173}{Reduction of the group of rationality by adjoining a rational function}{663}
\entry{174}{Reduction of the group to a pure Lie group $G_r$}{668}
\entry{175}{Necessary conditions for integrability by quadratures}{672}
\entry{176}{General linear differential equation. -- Simplicity of the group $L_1$}{677}
\entry{177}{The second part of Vessiot's theorem}{684}
\addlinespace[0.5em]
\entry{}{\textsc{Note I.} The finite transformations of a group $G_r=[X_1f,X_2f,\ldots,X_rf]$ composed with $r$ transformations taken respectively in the $r$ subgroups $[X_1f]$, $[X_2f]$, $\ldots$, $[X_rf]$}{688}
\entry{}{\textsc{Note II.} Ordinary differential equations possessing systems of fundamental solutions (Lie systems)}{692}
\bottomrule
\end{longtable}
\normalsize



\mainmatter
\chapter*{Chapter I}
\addcontentsline{toc}{chapter}{Chapter I}

\begin{center}
\begin{minipage}{0.88\linewidth}
\centering
Transformations and their composition -- Series of transformations and essential parameters -- Concept of a finite continuous group according to Lie -- Fundamental differential equations.\par
Groups with one parameter -- Infinitesimal transformations -- The $\infty^r$ independent infinitesimal transformations for a group with $r$ parameters.
\end{minipage}
\end{center}

\section*{\S{} 1. Transformations and their composition}

Let $x_1,x_2,\ldots,x_n$ be $n$ independent variables, which in general we shall suppose to be complex. We shall call a \emph{transformation} the passage from these $n$ variables to $n$ new independent variables $x'_1,x'_2,\ldots,x'_n$, connected with the primitive variables by a system of formulae:
\begin{equation*}\tag{1}
\left\{
\begin{aligned}
 x'_1&=f_1(x_1,x_2,\ldots,x_n),\\
 x'_2&=f_2(x_1,x_2,\ldots,x_n),\\
 &\hspace{1.6em}\cdots\cdots\cdots\\
 x'_n&=f_n(x_1,x_2,\ldots,x_n),
\end{aligned}
\right.
\end{equation*}
or, as we shall write more briefly,
\begin{equation*}\tag{$1^*$}
 x'_i=f_i(x_1,x_2,\ldots,x_n)\qquad (i=1,2,\ldots,n).
\end{equation*}
With respect to the functions $f_1,f_2,\ldots,f_n$ we shall always suppose, unless the contrary is stated explicitly, that they are analytic functions of the variables $x$, defined therefore in a primitive domain by an absolutely convergent series proceeding according to integral positive powers of the arguments
\[
 x_1-x_1^{(0)},\quad x_2-x_2^{(0)},\quad \ldots,\quad x_n-x_n^{(0)},
\]
where by $x_i^{(0)}$ we denote fixed values of the $x$. The functions $f_i$, thus defined at first only in the neighbourhood of a point $(x_1^{(0)},\ldots,x_n^{(0)})$, may then be analytically continued according to the principles of the general theory of analytic functions.

The hypothesis that the $x'_i$ are independent of one another entails that the functional determinant
\[
 \frac{\partial(f_1,f_2,\ldots,f_n)}{\partial(x_1,x_2,\ldots,x_n)}
\]
is in general different from zero. It is then well known that, by suitably restricting the neighbourhood of the point $x_i^{(0)}$, one can solve (1) with respect to the $x$, thus expressing the $x$ in terms of the $x'$ by the formulae of the inverse transformation
\begin{equation*}\tag{2}
 x_i=F_i(x'_1,x'_2,\ldots,x'_n).
\end{equation*}
As has been said, in general the variables will be regarded as complex. From the development itself of the investigations, however, it will appear that they are still applicable to real variables and transformations, provided one supposes that the functions $f_i$ on which the transformation depends are single-valued, finite, and continuous in a certain domain, together with their first, second, and so on derivatives, up to the order that it is necessary to consider.

A transformation such as (1) will often be denoted by a single letter, or symbol, $S$; in that case the inverse transformation (2) will be denoted by $S^{-1}$. What is essential in the transformation $S$ represented by (1) is naturally the species of the functions $f_i$, and not the way in which the primitive variables and the transformed variables are denoted; and so the transformation
\[
 y'_i=f_i(y_1,y_2,\ldots,y_n)
\]
will still be the same $S$.

Consider a first transformation
\[
 S)\qquad x'_i=f_i(x_1,x_2,\ldots,x_n)
\]
and a second
\[
 T)\qquad x''_i=\varphi_i(x'_1,x'_2,\ldots,x'_n).
\]
If we suppose that the variables $x$ remain in such a domain that the resulting values of the $x'$ belong to the domain in which the $\varphi_i$ are defined, we may perform the two transformations successively, and the transformation that results,
\[
 x''_i=\varphi_i\{f_1(x_1,\ldots,x_n),f_2(x_1,\ldots,x_n),\ldots,f_n(x_1,\ldots,x_n)\},
\]
or, as we shall often write for brevity,
\[
 x''_i=\varphi_i(f_1,f_2,\ldots,f_n),
\]
is called the composite transformation, or the product of the two transformations in the established order. It will be denoted by the symbol $ST$, placing on the left the transformation performed first.

Observe that by composing a transformation $S$ with its inverse $S^{-1}$, the result is the identical substitution, or the identity,
\[
 x'_i=x_i,
\]
which is denoted by the symbol $1$, so that
\[
 SS^{-1}=1.
\]
With the definition of the product of two transformations there is likewise established the meaning of the product of several transformations $S_1,S_2,\ldots,S_m$, taking
\[
 S_1S_2\cdots S_m=(S_1S_2\cdots S_{m-1})S_m.
\]
One should at once observe that in a product of several transformations the commutative law does not hold in general, but the associative law does hold. This property, which follows immediately from the very meaning of transformations,\footnote{It is enough to prove that the elementary associative law expressed by the formula $(S_1S_2)S_3=S_1(S_2S_3)$ holds. Now if $S_1$ carries the $x_i$ into the $x'_i$, $S_2$ carries the $x'_i$ into the $x''_i$, and finally $S_3$ carries the $x''_i$ into the $x'''_i$, the effect of the transformation on the left is to carry the $x_i$ into the $x'''_i$. On the other hand, since $S_2S_3$ carries the $x'_i$ into the $x'''_i$, that on the right also transforms the $x_i$ into the $x'''_i$, and therefore the two transformations are identical.} is important because it allows one to perform on a product of transformations the reductions that, by the associative law, are allowed in a product of ordinary quantities; thus, for example,
\[
 (S_1S_2)\cdot(S_3S_4)=S_1S_2S_3S_4,
\]
\[
 (S_1S)\cdot(S^{-1}S_2)=S_1S_2,\quad\hbox{etc.}
\]

When, on the old variables $x$ and the new variables $x'$ in a transformation
\[
 S)\qquad x'_i=f_i(x_1,x_2,\ldots,x_n)
\]
one performs the same transformation
\[
T\left\{
\begin{aligned}
 y_i&=\varphi_i(x_1,x_2,\ldots,x_n),\\
 y'_i&=\varphi_i(x'_1,x'_2,\ldots,x'_n),
\end{aligned}
\right.
\]
the transformation that arises in passing from the $y$ to the $y'$ is precisely
\[
 T^{-1}ST,
\]
which is called the transform of $S$ by means of $T$. Indeed, with $T^{-1}$ one passes from the $y$ to the $x$, with $S$ from the $x$ to the $x'$, and finally with $T$ from the $x'$ to the $y'$; hence, with $T^{-1}ST$, from the $y$ to the $y'$.

\section*{\S{} 2. Series $\infty^r$ of transformations}

Suppose that in the expressions of a transformation $S$ there enter $r$ parameters $a_1,a_2,\ldots,a_r$, so that one has
\begin{equation*}\tag{3}
 x'_i=f_i(x_1,x_2,\ldots,x_n; a_1,a_2,\ldots,a_r).
\end{equation*}
These parameters, like the variables $x$, will in general be supposed complex, and the $f_i$ will be analytic functions also of the parameters $a$.

We shall say that the $r$ parameters $a$ are essential in the $f$ when these functions cannot be reduced to depending on a smaller number of parameters $b$ which are functions of the $a$.\footnote{If, for example, the two parameters $a_1,a_2$ occur in the $f$ only in the combination $a_1+a_2$, the number of effective parameters would evidently be lowered at least to $r-1$.}
Letting the parameters $a$ run through all the values of a certain domain, the transformation $S$ will run through a continuous series of transformations. But it should be noted that, when the parameters $a_1,a_2,\ldots,a_r$ are not all essential, each transformation of the series will be obtained infinitely many times from different particular systems of parameter values, whereas, when they are essential, to that transformation there will correspond a single system of values of the parameters $a_1,a_2,\ldots,a_r$, or else a discrete number of such systems.

The first assertion is immediately evident, and the second is proved by observing that, in the series expansions that define the $f_i$ as functions of the $x$, the coefficients $\alpha$ are certain functions of the $a$, and the number $s$ of these independent $\alpha$ is certainly $\leq r$. When $s<r$, introducing the parameters $\alpha$ in place of the parameters $a$ shows that the latter are not essential; and conversely, if the parameters $a$ are not essential, then $s<r$, since the $a$ can be written as functions of those parameters $b$, fewer in number than $r$, which can be substituted for the $a$.

When the parameters $a$ are essential, and therefore $s=r$, and the equalities
\[
 f_i(x,a)=f_i(x,a')
\]
require that one have
\[
\begin{aligned}
 \alpha_1(a)&=\alpha'_1(a'),\\
 \alpha_2(a)&=\alpha'_2(a'),\\
 &\cdots\cdots\cdots\\
 \alpha_r(a)&=\alpha'_r(a'),
\end{aligned}
\]
and therefore the admissible values for $a'_1,\ldots,a'_r$ either reduce to the single system $a_1,a_2,\ldots,a_r$ or form a discrete series, as was to be proved.

From the circumstances just observed, whenever in (3) the parameters $a_1,a_2,\ldots,a_r$ are essential, one will say that the series (3) is a series $\infty^r$ of transformations. Strictly, when the parameters $a$ are complex the series is $\infty^{2r}$, but the locution introduced and now in use in similar cases cannot generate confusion if one remembers that, in the case of complex parameters, the effective degree of infinity of the series is always the double of that adopted.

It is important to establish a criterion for recognizing, given a series (3) of transformations, whether the $r$ parameters $a$ are essential or not. For this purpose we begin by observing that, if $a_1,a_2,\ldots,a_r$ are not essential, but can be reduced to the $s<r$ parameters
\[
 b_1,b_2,\ldots,b_s,
\]
which are functions of the $a$, we can determine such $r$ functions $A$ of the $a$ that the $n$ functions $f_i$, considered as functions of the $a$, all satisfy the homogeneous linear partial differential equation
\begin{equation*}\tag{$\alpha$}
 A_1\frac{\partial\psi}{\partial a_1}+A_2\frac{\partial\psi}{\partial a_2}+\cdots+A_r\frac{\partial\psi}{\partial a_r}=0.
\end{equation*}
It is enough, in fact, to determine the $A$ so that they satisfy the $s<r$ homogeneous linear equations
\begin{equation*}\tag{$\beta$}
 A_1\frac{\partial b_i}{\partial a_1}+A_2\frac{\partial b_i}{\partial a_2}+\cdots+A_r\frac{\partial b_i}{\partial a_r}=0\qquad (i=1,2,\ldots,s),
\end{equation*}
and then, since the $a$ enter in the $f_i$ only in the combinations $b_1,\ldots,b_s$, it is seen at once that all the $f$ satisfy $(\alpha)$.

One should further observe that the $s$ equations $(\beta)$, homogeneous and linear in the $A$, have characteristic precisely $=s$, because $b_1,b_2,\ldots,b_s$ are supposed independent (otherwise the number of essential parameters would still be $<s$); hence $r-s$ of the $A$ remain arbitrary, that is, in place of the single equation $(\alpha)$ we can write a system of $r-s$ such independent equations, which will form a complete system, having the $r-(r-s)$ solutions $b$, and of this system $(\alpha)$ the functions $f_1,f_2,\ldots,f_n$ will be just as many solutions.

Conversely, if there exists an equation $(\alpha)$, or a system of such equations, with the $A$ functions only of the $a$, of which the $f_i$ are solutions, then in these functions the parameters $a_1,a_2,\ldots,a_r$ will not all be essential, but will reduce to $s$ only, provided that there exist $r-s$ independent equations, such as $(\alpha)$, of which the $f_i$ are solutions. Indeed, in that hypothesis the system of the $r-s$ equations $(\alpha)$, supposed complete, will admit only $s=r-(r-s)$ independent solutions
\[
 b_1,b_2,\ldots,b_s,
\]
and the $f_i$, therefore, insofar as they depend on the $a$ and are by hypothesis solutions of the system $(\alpha)$, will reduce to functions of the $b$ alone.

We conclude from this. A necessary and sufficient condition that, in the $n$ functions
\[
 f_i(x_1,x_2,\ldots,x_n; a_1,a_2,\ldots,a_r),
\]
the $r$ parameters $a$ should not be essential is that there exist one or more equations of the form
\begin{equation*}\tag{$\alpha$}
 A_1\frac{\partial f}{\partial a_1}+A_2\frac{\partial f}{\partial a_2}+\cdots+A_r\frac{\partial f}{\partial a_r}=0
\end{equation*}
with the $A$ functions of the $a$, of which all the $f_i$ are solutions. Precisely, if there exist $r-s$ independent equations $(\alpha)$, forming a complete system of which the $f_i$ are solutions, the essential parameters in the $f$ reduce to $s$ only.

\section*{\S{} 3. Criterion for recognizing the number of essential parameters}

It remains to transform the theoretical criterion just established into another one of practical application. But first we observe that the considerations of the preceding \S{} apply not only to the case of a number $n$ of functions $f_i$ equal to the number of independent variables $x_1,x_2,\ldots,x_n$, but more generally to recognizing, given an arbitrary number $m$ of functions
\[
 f_i(x_1,x_2,\ldots,x_n;a_1,a_2,\ldots,a_r)\qquad (i=1,2,\ldots,m),
\]
of the $n$ independent variables $x$ and of $r$ parameters $a$, to how many essential parameters the parameters $a$ can be reduced.

Everything is reduced to finding the number $r-s$ of independent linear equations
\[
 A_1\frac{\partial f}{\partial a_1}+A_2\frac{\partial f}{\partial a_2}+\cdots+A_r\frac{\partial f}{\partial a_r}=0,
\]
with the $A$ functions of the $a$ alone, which are identically satisfied by all the functions $f_1,f_2,\ldots,f_m$. Since, by hypothesis, the $m$ equations
\begin{equation*}\tag{4}
 A_1\frac{\partial f_i}{\partial a_1}+A_2\frac{\partial f_i}{\partial a_2}+\cdots+A_r\frac{\partial f_i}{\partial a_r}=0\qquad (i=1,2,\ldots,m)
\end{equation*}
must hold, we begin by constructing the matrix
\[
M=\left\|\begin{array}{cccc}
\dfrac{\partial f_1}{\partial a_1} & \dfrac{\partial f_2}{\partial a_1} & \cdots & \dfrac{\partial f_m}{\partial a_1}\\[0.65em]
\dfrac{\partial f_1}{\partial a_2} & \dfrac{\partial f_2}{\partial a_2} & \cdots & \dfrac{\partial f_m}{\partial a_2}\\
\vdots & \vdots & & \vdots\\
\dfrac{\partial f_1}{\partial a_r} & \dfrac{\partial f_2}{\partial a_r} & \cdots & \dfrac{\partial f_m}{\partial a_r}
\end{array}\right\|
\]
and let $\mu$ be its characteristic. Certainly $\mu\le r$; but if $\mu=r$, then in (4) all the $A$ necessarily vanish, that is, the $r$ parameters $a$ are all essential and the question is settled. The number denoted above by $s$ is then $=r$.

If, differentiating the (4) once with respect to $x_1,x_2,\ldots,x_n$, one forms the new $nm$ homogeneous linear equations in the $A$
\begin{equation*}\tag{$4_1$}
 A_1\frac{\partial^2 f_i}{\partial a_1\partial x_k}+A_2\frac{\partial^2 f_i}{\partial a_2\partial x_k}+\cdots+A_r\frac{\partial^2 f_i}{\partial a_r\partial x_k}=0,
 \qquad\begin{gathered} i=1,2,\ldots,m,\\ k=1,2,\ldots,n,\end{gathered}
\end{equation*}
and constructs the complete matrix of the linear equations (4), $(4_1)$ in the $A$, which we shall denote by $M_1$,
\[
M_1=\left\|\begin{array}{ccccccc}
\dfrac{\partial f_1}{\partial a_1}&\dfrac{\partial f_2}{\partial a_1}&\cdots&\dfrac{\partial f_m}{\partial a_1}&
\dfrac{\partial^2 f_1}{\partial a_1\partial x_1}&\cdots&\dfrac{\partial^2 f_m}{\partial a_1\partial x_n}\\[0.65em]
\dfrac{\partial f_1}{\partial a_2}&\dfrac{\partial f_2}{\partial a_2}&\cdots&\dfrac{\partial f_m}{\partial a_2}&
\dfrac{\partial^2 f_1}{\partial a_2\partial x_1}&\cdots&\dfrac{\partial^2 f_m}{\partial a_2\partial x_n}\\
\vdots&\vdots&&\vdots&\vdots&&\vdots\\
\dfrac{\partial f_1}{\partial a_r}&\dfrac{\partial f_2}{\partial a_r}&\cdots&\dfrac{\partial f_m}{\partial a_r}&
\dfrac{\partial^2 f_1}{\partial a_r\partial x_1}&\cdots&\dfrac{\partial^2 f_m}{\partial a_r\partial x_n}
\end{array}\right\|,
\]
and which is formed from $M$ by adding to $f_1,f_2,\ldots,f_m$, as new functions, all their first derivatives with respect to the $x$, then, if $\mu_1$ denotes the characteristic of the matrix $M_1$, one will evidently have
\[
 \mu\le \mu_1\le r.
\]
If $\mu_1=r$, the equations (4) together with $(4_1)$ show that the $A$ are all zero and that the $r$ parameters $a$ are all essential. Supposing still that $\mu_1<r$, differentiate the $(4_1)$ again with respect to all the $x$ and form the new $n^2m$ equations
\begin{equation*}\tag{$4_2$}
 A_1\frac{\partial^3 f_i}{\partial a_1\partial x_k\partial x_l}+A_2\frac{\partial^3 f_i}{\partial a_2\partial x_k\partial x_l}+\cdots+A_r\frac{\partial^3 f_i}{\partial a_r\partial x_k\partial x_l}=0.
\end{equation*}
Let $\mu_2$ denote the characteristic of the matrix $M_2$ formed by adding, on the right of the columns of $M_1$, those relative to the equations $(4_2)$, that is, by adjoining to
\[
 f_i,\qquad \frac{\partial f_i}{\partial x_k},
\]
also the second derivatives $\dfrac{\partial^2 f_i}{\partial x_k\partial x_l}$. One will have $\mu_1\le\mu_2\le r$, and, if $\mu_2=r$, the $r$ parameters will be essential. Continuing in this way, we shall construct the successive matrices
\[
 M,
 M_1,
 M_2,
 \ldots
\]
of respective characteristics
\[
 \mu\le\mu_1\le\mu_2\le\mu_3\le\cdots .
\]
Since all the $\mu$ are $\le r$, the series
\[
 \mu,\mu_1,\mu_2,\ldots
\]
cannot be always increasing, and at a certain point one $\mu$ will necessarily be equal to the preceding one; let us say
\[
 \mu_k=\mu_{k-1}.
\]
Then it is seen at once that all the subsequent $\mu$ will retain the same value $\mu_{k-1}$. Indeed, since $\mu_k=\mu_{k-1}$, the columns added to the matrix $M_{k-1}$ to form the matrix $M_k$ are homogeneous linear combinations of $\mu_{k-1}$ fixed columns in $M_{k-1}$; these relations, differentiated with respect to the $x$, show that the new columns in $M_{k+1}$ are likewise homogeneous linear combinations of those same $\mu_{k-1}$ columns of $M_{k-1}$; the same is said for $M_{k+2},M_{k+3},\ldots$, as follows successively by differentiation. Hence the series of numbers
\[
 \mu,\mu_1,\mu_2,\ldots,\mu_{k-1},\ldots
\]
is always increasing up to a maximum value $\mu_{k-1}$, after which all the $\mu$ equal $\mu_{k-1}$; thus it will be enough to calculate the $\mu$ up to $\mu_{k-1}$.

If we call this maximum value of the $\mu$ by $s=\mu_{k-1}$, then certainly $s\le r$; hence in every case, since at least $\mu=1$, and the series of numbers
\[
 \mu,\mu_1,\mu_2,\ldots,\mu_{k-1}
\]
increases at each step by at least one unit, one will also have $k\le r$, so that, by constructing at most $r$ of our matrices, one will arrive at the end of the calculation.

We now prove that this number $s$, which gives the maximum value of the characteristics of the matrices
\[
 M,M_1,M_2,\ldots,
\]
is precisely equal to the number of essential parameters in the $m$ functions $f_i(x,a)$. For this, observe that the equations
\[
 (4),(4_1),(4_2),\ldots,(4_{k-1})
\]
reduce precisely to $s$ independent equations in the $A$, while it is useless to take account of the succeeding $(4_k),(4_{k+1}),\ldots$, because they are consequences of the $s$ already written.

Thus, if $s=r$, all the $A$ are necessarily zero and, since there exists no equation
\[
 A_1\frac{\partial f}{\partial a_1}+\cdots+A_r\frac{\partial f}{\partial a_r}=0,
\]
of which the $f_i$ are solutions, the $r$ parameters $a$ are in fact essential.

If, on the other hand, $s<r$, it is certainly possible to satisfy the said equations
\[
 (4),(4_1),(4_2),\ldots,(4_{k-1})
\]
with values of the $A$ which are in general functions not only of the $a$, but also of the $x$ themselves. In doing this one can take $r-s$ of the $A$ arbitrarily. We now show that, if for these $r-s$ arbitrary $A$ one takes functions only of the parameters $a$, then the remaining $s$ $A$ also will contain only the $a$. To fix ideas, suppose that the $r-s$ arbitrary $A$ are $A_1,A_2,\ldots,A_{r-s}$ and that the said equations then give
\[
 A_{r-s+1},\ldots,A_r
\]
as homogeneous linear functions of the first; let us write, for example,
\begin{equation*}\tag{5}
 A_{r-s+i}=\lambda_{i1}A_1+\lambda_{i2}A_2+\cdots+\lambda_{i,r-s}A_{r-s},
\end{equation*}
where the $\lambda$ are known functions of the $x$ and of the $a$. If we differentiate
\[
 (4),(4_1),\ldots,(4_{k-1})
\]
with respect to any $x$, say $x_\rho$, and observe that the $A$ also satisfy the equations $(4_1)$, we infer that the said equations are still satisfied upon replacing $A_1,A_2,\ldots,A_r$ by their derivatives
\[
 \frac{\partial A_1}{\partial x_\rho},\quad
 \frac{\partial A_2}{\partial x_\rho},\quad\ldots,\quad
 \frac{\partial A_r}{\partial x_\rho}
\]
with respect to one and the same $x$. Therefore, together with (5), the other relation also holds:
\[
 \frac{\partial A_{r-s+i}}{\partial x_\rho}=\lambda_{i1}\frac{\partial A_1}{\partial x_\rho}+\lambda_{i2}\frac{\partial A_2}{\partial x_\rho}+\cdots+\lambda_{i,r-s}\frac{\partial A_{r-s}}{\partial x_\rho};
\]
and it follows precisely that, taking $A_1,A_2,\ldots,A_{r-s}$ to be functions only of the parameters $a$, the $A_{r-s+1},\ldots,A_r$ also will contain only these parameters. Thus we have obtained exactly $r-s$ independent equations of the type
\[
 A_1\frac{\partial f}{\partial a_1}+A_2\frac{\partial f}{\partial a_2}+\cdots+A_r\frac{\partial f}{\partial a_r}=0,\tag{*}
\]
with the $A$ functions only of the parameters $a$, of which all the $f_i(x;a)$ are solutions; and therefore (\S{}2) the number of essential parameters is lowered precisely to $s$.\footnote{To obtain these $r-s$ independent equations it will be enough, for example, to put successively $A_1=1,A_2=0,A_3=0,\ldots,A_{r-s}=0$; $A_1=0,A_2=1,A_3=0,\ldots,A_{r-s}=0$; and each time to calculate the remaining $A$ from (5).}

We summarize the results obtained in the following rule. To see how many essential parameters $a$ there are in $m$ given functions
\[
 f_i(x_1,x_2,\ldots,x_n;a_1,a_2,\ldots,a_r)\qquad (i=1,2,\ldots,m)
\]
of the $x$ and of the parameters $a$, one begins by constructing the functional matrix $M$ of the derivatives of the $f$ with respect to the $a$; then the analogous matrix $M_1$, by adding to the $f_i$ all their first derivatives with respect to the $x$; next the matrix $M_2$, obtained by adding again the second derivatives, and so on. One prolongs the construction of these matrices $M,M_1,M_2,\ldots$ until the series $\mu,\mu_1,\mu_2,\ldots$ of their characteristics ceases to grow, after which the succeeding $\mu$ retain a constant value $s\le r$. This number $s$ gives precisely the number of essential parameters.

\section*{\S{} 4. Definition of finite continuous groups of transformations}

Transformations, being operations susceptible of composition, are subject to the general concepts of groups of operations. We shall say that a finite or infinite series of transformations constitutes a group when the product of any two transformations of the series, whether distinct or even coincident, is again contained in the same series.

Those groups which consist of a discrete number, finite or infinite, of transformations are called discontinuous. Such are, for example, the groups of substitutions on letters, whose theory has become, especially through the work of Galois, of fundamental importance for the study of algebraic equations. Other discontinuous, but infinite, groups occur in important branches of analysis: modular elliptic functions and automorphic functions.

The groups with which we shall be concerned in the present course are instead continuous groups, in whose transformations there enter arbitrary parameters capable of assuming continuously all the values of a certain domain. The theory of these continuous groups, as we possess it today, was built almost entirely by Sophus Lie; it has acquired great importance in mathematical science, both in the analytic field and in the geometric field, in particular for analysis in the theory of the integration of ordinary and partial differential equations, where the role of the theory of continuous groups is comparable to that which the theory of substitution groups plays in the study of algebraic equations.

According to Lie, continuous groups are called finite when their transformations depend on a finite number of arbitrary parameters. The simplest finite continuous groups are those whose transformations are defined by a single system of equations
\begin{equation*}\tag{6}
 x'_i=f_i(x_1,x_2,\ldots,x_n;a_1,a_2,\ldots,a_r)\qquad (i=1,2,\ldots,n),
\end{equation*}
with $r$ parameters, which may at once be supposed essential. The property of the series $\infty^r$ of transformations (6) of constituting a group is translated into the following functional property for the $f_i$. Since two successive transformations of the series (6)
\[
 x'_i=f_i(x_1,x_2,\ldots,x_n;a_1,a_2,\ldots,a_r),
\]
\[
 x''_i=f_i(x'_1,x'_2,\ldots,x'_n;b_1,b_2,\ldots,b_r)
\]
must compose into a third transformation
\[
 x''_i=f_i(x_1,x_2,\ldots,x_n;c_1,c_2,\ldots,c_r),
\]
the parameters $c$ of the composed transformation depending only on those $a,b$ of the components, the functions $f_i$ must satisfy the functional relations
\[
 f_i\{f_1(x,a),f_2(x,a),\ldots,f_n(x,a);b_1,b_2,\ldots,b_r\}=f_i(x_1,x_2,\ldots,x_n;c_1,c_2,\ldots,c_r),
\]
where the $c$ are functions of the $a,b$, say
\[
 c_i=\varphi_i(a_1,a_2,\ldots,a_r;b_1,b_2,\ldots,b_r).
\]

In general, when we speak of a group of transformations, we shall understand that it is a finite continuous group whose transformations are contained in a single series (6). The other, more general finite groups which it is useful to consider are those in which the transformations of the group, instead of constituting a single series (6), constitute several such series and even an infinite but discrete number of them. Moreover the study of such more general finite groups, which are called mixed groups, is reduced, as we shall prove in the next chapter, to that of the simple groups (6) combined with ordinary discontinuous groups.

Every continuous group that is not finite could be called an infinite continuous group. However, this denomination is adopted by Lie only in the case, to which we shall return later (see \S{}36), in which the transformations of the group itself are defined by a certain system of partial differential equations. For the rest, in this course we shall not concern ourselves at all with infinite groups in Lie's sense or in a more general sense, and we shall always understand that the group is finite.

An observation is immediately necessary in order to avoid an easy error. In ordinary groups of a finite number of operations there always exists in the group the identity operation, as well as, together with each operation, its inverse. But for continuous groups we make the observation, which also holds for infinite discontinuous groups, that from the definition of a continuous group it does not at all follow that the group must always contain the identity transformation, and still less that together with each transformation of the group it must contain its inverse.

We shall prove later only that, if the group contains the identity transformation, corresponding to the values
\[
 a_1^{(0)},a_2^{(0)},\ldots,a_r^{(0)}
\]
of the parameters, and if these values belong to the interior of the domain where the $f_i$ are defined, then, when the $a$ remain in a sufficiently restricted neighborhood of $(a^{(0)})$, together with each transformation $S_a$ in the group there will effectively exist its inverse $S_a^{-1}$.

\section*{\S{} 5. Examples of finite continuous groups}

To illustrate the general concepts laid down in the preceding \S{}, we shall give some simple examples of continuous groups taken from geometry, where they present themselves in the most intuitive way.

$1^{\circ}$. If a figure moves in space without changing shape, the coordinates $x',y',z'$ of each point of the figure in the new position are expressed by integral linear functions of the primitive coordinates $x,y,z$. With orthogonal axes assumed, the formulae are written
\[
\left\{
\begin{aligned}
 x'&=a_1x+a_2y+a_3z+a,\\
 y'&=b_1x+b_2y+b_3z+b,\\
 z'&=c_1x+c_2y+c_3z+c,
\end{aligned}\right.\tag{a}
\]
where the nine constants
\[
\begin{vmatrix}
 a_1&a_2&a_3\\
 b_1&b_2&b_3\\
 c_1&c_2&c_3
\end{vmatrix}
\]
are the coefficients of a dextrorse orthogonal substitution and therefore are expressed in terms of only three independent quantities, for example the Euler angles $\theta,\varphi,\psi$. The totality of rigid motions of space, represented by the formulae (a), where
\[
\begin{vmatrix}
 a_1&a_2&a_3\\
 b_1&b_2&b_3\\
 c_1&c_2&c_3
\end{vmatrix}
\]
runs through all dextrorse orthogonal substitutions and $a,b,c$ assume all possible values, therefore constitutes a finite continuous group with 6 essential parameters. It is clear that in this group the identity transformation exists and, together with each transformation, its inverse.

$2^{\circ}$. If to all the motions (a) we adjoin the symmetry, for example with respect to the plane $xy$, given by the formulae
\[
 x'=x,
 \qquad y'=y,
 \qquad z'=-z,
\]
we obtain a second series of transformations which differs from (a) only in this, that the corresponding orthogonal substitutions
\[
\begin{vmatrix}
 a_1&a_2&a_3\\
 b_1&b_2&b_3\\
 c_1&c_2&c_3
\end{vmatrix}
\]
are sinistrorse. The two series together furnish a simple example of a mixed group.

$3^{\circ}$. The motions of plane figures in their own plane form a continuous group with three parameters; together with reflections they give a mixed group. The first is analytically represented by the formulae
\[
\left\{
\begin{aligned}
 x'&=x\cos\alpha-y\sen\alpha+a,\\
 y'&=x\sen\alpha+y\cos\alpha+b,
\end{aligned}\right.
\]
the second by adjoining to these the others
\[
\left\{
\begin{aligned}
 x'&=x\cos\alpha+y\sen\alpha+a,\\
 y'&=x\sen\alpha-y\cos\alpha+b.
\end{aligned}\right.
\]

$4^{\circ}$. If to the total group of motions of space we adjoin a similarity
\[
 x'=kx,
 \qquad y'=ky,
 \qquad z'=kz
\]
with variable parameter $k$, we obtain a continuous group with 7 parameters, the group of similarities.

$5^{\circ}$. By combining the motions of space with all inversions by reciprocal radius-vectors with respect to the sphere of space, one obtains a continuous group with 10 parameters, which is the group of conformal transformations. According to a theorem of Liouville, it contains all and only the transformations of space that preserve angles.

$6^{\circ}$. The totality of collineations of space forms a continuous group with 15 parameters, the so-called projective group of $S_3$. In it, those collineations that leave fixed a non-degenerate conic form a subgroup with 4 parameters; those that leave fixed a non-degenerate quadric form a subgroup with 6 parameters; and finally those that carry into itself a skew cubic form a subgroup with 3 parameters.

Similarly, those collineations of space that leave fixed a non-special linear complex form a group with 10 parameters. In all the examples given above the transformation formulae are given by linear substitutions, except in example $5^{\circ}$, where we have quadratic substitutions.\footnote{This is a particular example of the Cremonian groups, in which the transformation formulae give the new variables expressed rationally in terms of the old ones and conversely.}

The generalization to space of $n$ dimensions is immediate. Thus the homogeneous linear substitutions on $n$ variables
\begin{equation*}\tag{$\beta$}
 x'_i=\sum\limits_{k=1}^n a_{ik}x_k,
\end{equation*}
with the $n^2$ parameters $a_{ik}$, evidently give us a group of Lie. Composing the preceding one with
\[
 x''_i=\sum\limits_{k=1}^n b_{ik}x'_k
\]
gives the composed substitution
\[
 x''_i=\sum\limits_{k=1}^n c_{ik}x_k,
\]
where the $c$ are expressed in terms of the $a,b$ by the formulae
\[
 c_{hk}=\sum\limits_i^{1,\ldots,n} a_{hi}b_{ik}.
\]
If in $(\beta)$ we impose among the parameters $a$ the relation expressed by
\[
 a=\begin{vmatrix}
 a_{11}&a_{12}&\cdots&a_{1n}\\
 \cdots&\cdots&\cdots&\cdots\\
 a_{n1}&a_{n2}&\cdots&a_{nn}
\end{vmatrix}=1,
\]
we shall evidently have a subgroup with $n^2-1$ parameters. If, instead, we make the modulus $a$ of the substitution $(\beta)$ run only through the $\nu$ $\nu$-th roots of unity, $\nu$ being a fixed integer, we shall evidently have a mixed group with $\nu$ series of transformations with $n^2-1$ parameters. Likewise the dextrorse orthogonal substitutions $(\beta)$ form a group with $\dfrac{n(n-1)}{2}$ parameters.

\section*{\S{} 6. Solubility of the equations $c_i=\varphi_i(a,b)$ with respect to the $a$ or the $b$}

Suppose that we have a series $\infty^r$ of transformations
\begin{equation*}\tag{6}
 x'_i=f_i(x_1,x_2,\ldots,x_n;a_1,a_2,\ldots,a_r)\qquad (i=1,2,\ldots,n),
\end{equation*}
which form a group. From the functional equations
\begin{equation*}\tag{7}
 f_i(x',b)=f_i(x,c),
\end{equation*}
which are identically satisfied when, for the $x'$, one puts their values as functions of the $x,a$,
\begin{equation*}\tag{$7^*$}
 x'_i=f_i(x,a)
\end{equation*}
and, for the $c$, their values as functions of the $a,b$,
\begin{equation*}\tag{8}
 c_i=\varphi_i(a,b),
\end{equation*}
one can draw some important preliminary consequences.

We first prove that the (8) are soluble with respect to the $b$, that is, that the functional determinant
\[
 \frac{\partial(\varphi_1,\varphi_2,\ldots,\varphi_r)}{\partial(b_1,b_2,\ldots,b_r)}
\]
cannot be zero. In the contrary case, indeed, one could determine functions $\alpha_1,\alpha_2,\ldots,\alpha_r$ of the $a,b$ satisfying the system of the $r$ homogeneous linear equations
\begin{equation*}\tag{9}
 \sum\limits_{l=1}^r \alpha_l\frac{\partial \varphi_i}{\partial b_l}=0\qquad (i=1,2,\ldots,r).
\end{equation*}
Now, differentiating the identity (7), where the substitutions $(7^*),(8)$ are imagined to have been carried out, with respect to $b_k$, we obtain
\[
 \frac{\partial f_i(x',b)}{\partial b_k}
 =\sum\limits_{j=1}^r\frac{\partial f_i(x,c)}{\partial c_j}\frac{\partial\varphi_j}{\partial b_k};
\]
hence, multiplying by $\alpha_k$ and summing from $k=1$ to $k=r$, we infer from (9) that
\[
 \sum\limits_k^{1,\ldots,r} \alpha_k\frac{\partial f_i(x',b)}{\partial b_k}=0.
\]
If in this identity one were to put for the $a$ the values of any particular system that does not annul all the $\alpha$, these would become functions $B_k$ of the $b$ alone. The functions
\[
 f_i(x',b)
\]
would therefore be solutions of an equation of the type
\[
 B_1\frac{\partial f}{\partial b_1}+B_2\frac{\partial f}{\partial b_2}+\cdots+B_r\frac{\partial f}{\partial b_r}=0,
\]
whence one would conclude (\S{}2) that in the functions $f_i(x';b)$ the $r$ parameters $b$ would not all be essential, contrary to the supposition.

To prove that the relations (8) are also soluble with respect to the $a$, it is enough to observe that the following theorem holds. If the $\infty^r$ transformations (6) constitute a group, then their inverse transformations
\[
 x_i=F_i(x'_1,\ldots,x'_n;a_1,\ldots,a_r)
\]
also constitute a group with $r$ parameters. Indeed, if we denote symbolically by $S_a$ the transformation (6) corresponding to the values $a$ of the parameters, since the (6) form a group,
\[
 S_aS_b=S_c,
\]
and therefore
\[
 S_b^{-1}S_a^{-1}=S_c^{-1},
\]
which proves that the inverse transformations also form a group. The number of parameters in the inverse group is again $=r$, since, in passing from a group to its inverse, the number of parameters evidently cannot increase and hence cannot decrease either.\footnote{Analytically one can establish this result by using the identities
\[
 f_i(F_1(x,a),\ldots,F_n(x,a);a_1,a_2,\ldots,a_r)=x_i,
\]
which, differentiated with respect to $a_k$, give
\[
 \frac{\partial F_i}{\partial a_k}+\sum\limits_{h=1}^n\frac{\partial f_i}{\partial x_h}\frac{\partial F_h}{\partial a_k}=0.
\]
If $a_1,a_2,\ldots,a_r$ were not all essential in the $F$, these would all satisfy an equation of the form
\[
 A_1\frac{\partial F}{\partial a_1}+\cdots+A_r\frac{\partial F}{\partial a_r}=0
\]
with the $A$ functions of the $a$. The preceding identities show that the $f_i$ also would satisfy it, and in these the parameters $a$ would therefore not all be essential.}
After this, since the relations (8) are the same ones which, for the primitive group, correspond to the equality $S_aS_b=S_c$ and, for the inverse group, to $S_b^{-1}S_a^{-1}=S_c^{-1}$, it is clear that the (8) are also soluble with respect to the $a$, as was to be proved.

We have therefore established the result. The equations (8)
\[
 c_i=\varphi_i(a,b),
\]
which, for a series $\infty^r$ of transformations constituting a group, connect the parameters $a,b$ of the component transformations $S_a,S_b$ with the parameters $c$ of the composite $S_c=S_aS_b$, are soluble both with respect to the $a$ and with respect to the $b$; that is, the functional determinants
\[
 \frac{\partial(\varphi_1,\varphi_2,\ldots,\varphi_r)}{\partial(a_1,a_2,\ldots,a_r)},
 \qquad
 \frac{\partial(\varphi_1,\varphi_2,\ldots,\varphi_r)}{\partial(b_1,b_2,\ldots,b_r)}
\]
are each different from zero.


\section*{\S{} 7. Fundamental differential equations}

We now proceed to a fundamental investigation for the theory of finite continuous groups, namely to prove how the functional equations (7) are translated, for the functions of transformation $f_i$ which occur in a series $\infty^r$ (6) constituting a group, into a system of total differential equations, unlimitedly integrable, which the $f_i$ themselves must satisfy when considered as functions of the parameters $a$.

The functional equations
\begin{equation*}\tag{7}
 f_i(x',b)=f_i(x,c)
\end{equation*}
must be translated, as we know, into identities with respect to the $x,a,b$, regarded as independent, when for the $x_i'$ and $c_i$ one puts the values
\begin{equation*}\tag{6}
 x_i'=f_i(x,a),
\end{equation*}
\begin{equation*}\tag{8}
 c_\alpha=\varphi_\alpha(a,b).
\end{equation*}
Since the (8) are soluble with respect to the $b$ (\S{}6), we may instead regard the (7) as identities in the $x,a,c$, where for the $x'$ one puts the values (6), and for the $b$ the values as functions of $a,c$ obtained from (8). Differentiating (7) with respect to any one of the $a$ under this connection, we obtain
\begin{equation*}\tag{10}
 \sum\limits_l^{1,\ldots,n}\frac{\partial f_i}{\partial x_l'}\frac{\partial x_l'}{\partial a_k}
 +\sum\limits_s^{1,\ldots,r}\frac{\partial f_i}{\partial b_s}\frac{\partial b_s}{\partial a_k}=0
 \qquad (i=1,2,\ldots,n).
\end{equation*}
Since the functional determinant
\[
 J=\frac{\partial(f_1,f_2,\ldots,f_n)}{\partial(x'_1,x'_2,\ldots,x'_n)}
\]
is not zero, we can solve the $n$ equations (10), linear in
\[
 \frac{\partial x'_1}{\partial a_k},\quad
 \frac{\partial x'_2}{\partial a_k},\quad\ldots,\quad
 \frac{\partial x'_n}{\partial a_k},
\]
with respect to these derivatives. If by $\Theta_{ij}$ we denote the quotient of the algebraic complement of $\dfrac{\partial f_j}{\partial x'_i}$ in the determinant $J$ by the determinant itself, the solution of (10) gives
\[
 \frac{\partial x'_i}{\partial a_k}
 =-\sum\limits_s^{1,\ldots,r}\sum\limits_j^{1,\ldots,n}
 \Theta_{ij}\frac{\partial f_j}{\partial b_s}\frac{\partial b_s}{\partial a_k}.
\]
Now the sum
\[
 -\sum\limits_j^{1,\ldots,n}\Theta_{ij}\frac{\partial f_j(x',b)}{\partial b_s}
\]
is a certain function of the $x',b$, which we denote by $\Phi_{si}(x',b)$, while the derivative $\dfrac{\partial b_s}{\partial a_k}$ is a function of $a,b$, which can be calculated from (8), and which we denote by
\[
 \frac{\partial b_s}{\partial a_k}=\psi_{sk}(a,b).
\]
With these notations the preceding formulae become
\begin{equation*}\tag{11}
 \frac{\partial x'_i}{\partial a_k}
 =\sum\limits_s^{1,\ldots,r}\psi_{sk}(a,b)\Phi_{si}(x',b).
\end{equation*}
These formulae, like the (7) from which they were obtained by differentiation, are identically satisfied when one puts
\[
 x'_i=f_i(x,a),
\]
since the $x,a,b$ are completely independent. In the first members of (11)
\[
 \frac{\partial f_i(x,a)}{\partial a_k}
\]
the $b$ do not occur; hence one sees that in the second members also they can be only apparent. If, therefore, we give the $b$ a fixed system of values, otherwise arbitrary,
\[
 b_i=b_i^{(0)},
\]
and put
\[
 \Phi_{sj}\bigl(x',b^{(0)}\bigr)=\xi_{sj}(x'_1,x'_2,\ldots,x'_n),
\]
\[
 \psi_{sk}\bigl(a,b^{(0)}\bigr)=\psi_{sk}(a_1,a_2,\ldots,a_r),
\]
the (11) give us the required system of fundamental equations:
\begin{equation*}\tag{A}
 \frac{\partial x'_i}{\partial a_k}
 =\sum\limits_s^{1,\ldots,r}\xi_{si}(x')\psi_{sk}(a)
 \qquad
 \begin{cases}
 i=1,2,\ldots,n,\\
 k=1,2,\ldots,r,
 \end{cases}
\end{equation*}
which is satisfied by the functions
\[
 x'_i=f_i(x_1,x_2,\ldots,x_n;a_1,a_2,\ldots,a_r)
 \qquad i=1,2,\ldots,n,
\]
defining our group.

Our equations (A) may also be written in the equivalent form
\begin{equation*}\tag{$A^*$}
 dx'_j=\sum\limits_k^{1,\ldots,r}\left(\sum\limits_s^{1,\ldots,r}\xi_{sj}(x')\psi_{sk}(a)\right)da_k
 \qquad (j=1,2,\ldots,n)
\end{equation*}
and constitute a system of total differential equations. In the integrals
\begin{equation*}\tag{$\gamma$}
 x'_j=f_j(x,a)
\end{equation*}
the $n$ variables $x$ occur as arbitrary constants; and since the $(\gamma)$ are soluble with respect to the $x$, it follows that our system of total differential equations (A), or $(A^*)$, is unlimitedly integrable.

We summarize the results obtained so far in the theorem:

If the equations
\[
 x'_i=f_i(x_1,\ldots,x_n; a_1,a_2,\ldots,a_r)
 \qquad i=1,2,\ldots,n
\]
represent a group with $r$ essential parameters $a$, then the functions $f_i(x,a)$, considered as depending on the $a$, satisfy a system of total differential equations of the form
\begin{equation*}\tag{A}
 \frac{\partial x'_i}{\partial a_k}
 =\sum\limits_s^{1,\ldots,r}\xi_{si}(x')\psi_{sk}(a)
 =\xi_{1i}(x')\psi_{1k}(a)+\xi_{2i}(x')\psi_{2k}(a)+\cdots+\xi_{ri}(x')\psi_{rk}(a),
\end{equation*}
\footnote{Observe the special form of the second members of these equations, as sums of $r$ products of functions of the $a$ alone by functions of the $x'$ alone.}
this system is unlimitedly integrable, and the $f_i(x,a)$ give its general integrals.

As an example take the group of motions in the plane (\S{}5), represented by the equations
\[
\begin{cases}
 x'_1=x_1\cos a_3-x_2\sen a_3+a_1,\\
 x'_2=x_1\sen a_3+x_2\cos a_3+a_2,
\end{cases}
\]
where $x_1,x_2$ represent orthogonal Cartesian coordinates and $a_1,a_2,a_3$ are parameters. The differential equations are here
\[
\begin{array}{lll}
 \dfrac{\partial x'_1}{\partial a_1}=1, &
 \dfrac{\partial x'_1}{\partial a_2}=0, &
 \dfrac{\partial x'_1}{\partial a_3}=a_2-x'_2,\\[0.9em]
 \dfrac{\partial x'_2}{\partial a_1}=0, &
 \dfrac{\partial x'_2}{\partial a_2}=1, &
 \dfrac{\partial x'_2}{\partial a_3}=x'_1-a_1,
\end{array}
\]
and they fall under the general form (A) by putting
\[
 \left\|\begin{array}{cc}
 \xi_{11}&\xi_{12}\\
 \xi_{21}&\xi_{22}\\
 \xi_{31}&\xi_{32}
 \end{array}\right\|
 =
 \left\|\begin{array}{cc}
 1&0\\
 0&1\\
 x'_2&-x'_1
 \end{array}\right\|,
 \qquad
 \left\|\begin{array}{ccc}
 \psi_{11}&\psi_{12}&\psi_{13}\\
 \psi_{21}&\psi_{22}&\psi_{23}\\
 \psi_{31}&\psi_{32}&\psi_{33}
 \end{array}\right\|
 =
 \left\|\begin{array}{ccc}
 1&0&a_2\\
 0&1&-a_1\\
 0&0&-1
 \end{array}\right\|.
\]


\section*{\S{} 8. Some properties of the functions $\psi_{sh}(a)$, $\xi_{si}(x')$}

We have established the result, of fundamental importance, that the functions $f_i(x,a)$ which give a series $\infty^r$ of transformations
\begin{equation*}\tag{6}
 x'_i=f_i(x,a)
\end{equation*}
constituting a group necessarily satisfy a system of differential equations of the form (A). We say at once, however, that the conditions expressed by (A), necessary in order that the transformations (6) constitute a group, are not altogether sufficient for this purpose,\footnote{An example at once shows the point. The simple system
\[
 \frac{dx'_1}{da}=\frac{x'_1}{x'_2},\qquad \frac{dx'_2}{da}=1
\]
with a single parameter $a$ has the typical form (A), by taking $\xi_{11}=x'_1/x'_2$, $\xi_{12}=1$, $\psi_{11}(a)=1$; but the general integrals
\[
 x'_1=x_1(x_2+a),\qquad x'_2=x_2+a
\]
do not give the transformations of a group.} but that certain complementary conditions are still required, and these will be examined in the next chapter. Meanwhile we observe some simple but important properties which hold for every series $\infty^r$ of transformations (6), when the differential equations of a system (A) are satisfied, even if that series should not constitute a group.

In the first place we say that the determinant of the functions $\psi_{hk}(a)$
\[
 \psi=\left|\begin{array}{cccc}
 \psi_{11}&\psi_{12}&\cdots&\psi_{1r}\\
 \psi_{21}&\psi_{22}&\cdots&\psi_{2r}\\
 \cdots&\cdots&\cdots&\cdots\\
 \psi_{r1}&\psi_{r2}&\cdots&\psi_{rr}
 \end{array}\right|
\]
is necessarily different from zero. In the contrary case, in fact, we could determine $r$ functions
\[
 A_1,A_2,\ldots,A_r
\]
of the $a$, not all zero, such that they satisfy the system of the $r$ homogeneous linear equations
\[
 \sum\limits_h^{1,\ldots,r} A_h\psi_{hs}(a)=0\qquad (s=1,2,\ldots,r).
\]
From the differential equations (A) there would then result the following equations
\[
 \sum\limits_h^{1,\ldots,r} A_h\frac{\partial f_i(x,a)}{\partial a_h}=0\qquad (i=1,2,\ldots,n),
\]
and therefore (\S{}2) the $r$ parameters $a$ would not be essential, contrary to the hypothesis.

After this, if we consider the $r$ equations (A) which are obtained by keeping $i$ fixed and putting $k=1,2,\ldots,r$, it is clear that we can solve them precisely with respect to
\[
 \xi_{1i},\xi_{2i},\ldots,\xi_{ri},
\]
and the resolving formulae will take the following form:
\begin{equation*}\tag{B}
 \xi_{ji}(x')=\sum\limits_k^{1,\ldots,r}\alpha_{jk}(a)\frac{\partial x'_i}{\partial a_k},
\end{equation*}
where $\alpha_{jk}$ denotes the algebraic complement of $\psi_{jk}$ in $\psi$, divided by $\psi$ itself.

It follows from this an important result for what follows, namely that it is impossible to determine $r$ constants, not all zero,
\[
 e_1,e_2,\ldots,e_r
\]
which satisfy the $n$ equations
\begin{equation*}\tag{12}
 \sum\limits_j^{1,\ldots,r}e_j\xi_{ji}(x')=0\qquad (i=1,2,\ldots,n).
\end{equation*}
Indeed, by (B), equations (12) would give us
\[
 \sum\limits_k^{1,\ldots,r}\left(\sum\limits_j^{1,\ldots,r}e_j\alpha_{jk}(a)\right)
 \frac{\partial f_i(x,a)}{\partial a_k}=0,
\]
and the parameters $a_1,a_2,\ldots,a_r$ would not all be essential.

We formulate the results obtained in the theorem:
If a series $\infty^r$ of transformations
\[
 x'_i=f_i(x,a)
\]
satisfies a system of differential equations of type (A), 1. the determinant of the $\psi$ is necessarily different from zero, and the equations (A), solved with respect to the $\xi$, can be put under the equivalent form (B); 2. it is impossible to determine $r$ constants $e_1,e_2,\ldots,e_r$ (non-zero) which satisfy the $n$ linear equations
\[
 \sum\limits_j^{1,\ldots,r}e_j\xi_{ji}(x')=0\qquad (i=1,2,\ldots,n).
\]

\section*{\S{} 9. Groups with one parameter}

We begin by applying these first principles of the theory to the case of groups with one parameter; these form, in a certain sense, the elements from which finite continuous groups (and also infinite groups) are composed, and one may say that in the general theory of continuous groups they perform an office analogous to that of groups of powers in the theory of discontinuous groups.

If the series $\infty^1$ of transformations
\begin{equation*}\tag{13}
 x'_i=f_i(x_1,x_2,\ldots,x_n;a)\qquad (i=1,2,\ldots,n)
\end{equation*}
with one parameter $a$ constitutes a group, then by the fundamental theorem of \S{}7 the $x'_i$, considered as functions of $a$, will satisfy a system of ordinary differential equations of the form
\begin{equation*}\tag{14}
 \frac{dx'_i}{da}=\psi(a)\cdot \xi_i(x'_1,x'_2,\ldots,x'_n),
 \qquad i=1,2,\ldots,n,
\end{equation*}
and will constitute its general integral.

To determine the $f_i(x,a)$ completely, it will therefore suffice to know the form which these $f_i$, that is, the $x'_i$, assume for a particular value of $a$, say for $a=a_0$.

It is clear that the integral system of (14), that is, of the equations
\[
 \frac{dx'_1}{\xi_1}=\frac{dx'_2}{\xi_2}=\cdots=\frac{dx'_n}{\xi_n}=\psi(a)\,da,
\]
can be put under the form solved with respect to the arbitrary constants
\begin{equation*}\tag{15}
\left\{
\begin{aligned}
 \Omega_1(x'_1,x'_2,\ldots,x'_n)&=C_1,\\
 \Omega_2(x'_1,x'_2,\ldots,x'_n)&=C_2,\\
 &\hspace{1.5em}\cdots\cdots\cdots\\
 \Omega_{n-1}(x'_1,x'_2,\ldots,x'_n)&=C_{n-1},\\
 \Omega_n(x'_1,x'_2,\ldots,x'_n)&=\int_{a_0}^{a}\psi(a)\,da+C_n,
\end{aligned}
\right.
\end{equation*}
where the $\Omega_i(x'_1,x'_2,\ldots,x'_n)$ are certain $n$ independent functions of $x'_1,x'_2,\ldots,x'_n$, and $C_1,C_2,\ldots,C_n$ are arbitrary constants.

Suppose now that in the series (13) the identity substitution is contained and corresponds to the value $a_0$ of $a$, so that one has
\[
 x'_i=x_i\qquad\text{for }a=a_0.
\]
Putting $a=a_0$ in (15), one obtains
\[
 C_i=\Omega_i(x_1,x_2,\ldots,x_n),
\]
and the integral system (15), changing for simplicity the parameter $a$ into
\[
 t=\int_{a_0}^{a}\psi(a)\,da,
\]
takes the form
\begin{equation*}\tag{16}
\left\{
\begin{aligned}
 \Omega_i(x'_1,x'_2,\ldots,x'_n)&=\Omega_i(x_1,x_2,\ldots,x_n)\qquad (i=1,2,\ldots,n-1),\\
 \Omega_n(x'_1,x'_2,\ldots,x'_n)&=\Omega_n(x_1,x_2,\ldots,x_n)+t,
\end{aligned}
\right.
\end{equation*}
that is, it is perfectly equivalent to (13) under the hypothesis
\[
 x'_i=x_i\quad\text{for }t=0\quad (a=a_0).
\]
The equations (16) display the group property in the simplest way, since, if we change the variables $x_i$ into
\[
 y_i=\Omega_i(x_1,x_2,\ldots,x_n)\qquad i=1,2,\ldots,n,
\]
they take the very simple form
\[
 y'_1=y_1,\quad y'_2=y_2,\quad\ldots,\quad y'_{n-1}=y_{n-1},\quad y'_n=y_n+t.
\]
This transformation will be called a translation of the space of $n$ dimensions $(y_1,y_2,\ldots,y_n)$.

In general, if in any group $G$ the same substitution $T$ is made simultaneously on the primitive variables and on the transformed variables, and also, if desired, the parameters are changed in an arbitrary manner, the considerations at the end of \S{}1 prove that the new transformations again run through a group, the group $T^{-1}GT$. Two such groups, obtained one from the other by a change of variables and of parameters, are called similar.

We may therefore state this first result:
If a continuous group with one parameter contains the identity, it is similar to a group of translations, and its transformations are mutually permutable and two by two inverse to one another.

Suppose now, more generally, that we have a series $\infty^1$ of transformations (13), which satisfies a system of differential equations of the form (14), and denote by
\[
 \varphi_1(x),\varphi_2(x),\ldots,\varphi_n(x)
\]
the functions of the $x$ to which the $f_i(x,a)$ reduce for $a=a_0$; the system (15) equivalent to (13) will be written:
\[
\left\{
\begin{aligned}
 \Omega_i(x'_1,x'_2,\ldots,x'_n)&=\Omega_i(\varphi_1,\varphi_2,\ldots,\varphi_n)
 \qquad (i=1,2,\ldots,n-1),\\
 \Omega_n(x'_1,x'_2,\ldots,x'_n)&=\Omega_n(\varphi_1,\varphi_2,\ldots,\varphi_n)+t.
\end{aligned}
\right.
\]
This transformation is obtained by carrying out, first, the transformation
\[
 y_1=\varphi_1(x_1,x_2,\ldots,x_n),\quad\ldots,\quad y_n=\varphi_n(x_1,x_2,\ldots,x_n),
\]
and then the other
\[
\left\{
\begin{aligned}
 \Omega_i(x'_1,x'_2,\ldots,x'_n)&=\Omega_i(y_1,y_2,\ldots,y_n)\qquad (i=1,2,\ldots,n-1),\\
 \Omega_n(x'_1,x'_2,\ldots,x'_n)&=\Omega_n(y_1,y_2,\ldots,y_n)+t.
\end{aligned}
\right.
\]
The first is a fixed transformation of the series (13) (corresponding to $a=a_0$), and the second, as $t$ varies, generates a group of the species considered above. We therefore have the theorem:

If a series $\infty^1$ of transformations (13) satisfies differential equations of type (14), its transformations are obtained by composing a fixed transformation of the series itself with all the transformations of the group with one parameter (with the identity substitution) defined by the same differential equations (14).

\section*{\S{} 10. Infinitesimal transformation of a $G_1$}

We return to the consideration of groups with one parameter containing the identity and therefore containing, together with every transformation, its inverse (\S{}9). Such groups will henceforth be denoted by the symbol $G_1$, the subscript referring to the number of parameters. The formulae relative to such a group $G_1$ are obtained, as has been seen, by integrating a system of differential equations
\[
 \frac{dx'_i}{dt}=\xi_i(x'_1,x'_2,\ldots,x'_n)
\]
with the initial conditions
\[
 x'_i=x_i\qquad\text{for }t=0.
\]
The functions $\xi_i(x')$ are supposed analytic; and, if one further assumes that the point $(x'_1,x'_2,\ldots,x'_n)$ is a regular point for these functions, we know that the integrals $x'_i$ will be expandable in power series in $t$ according to the formula
\[
 x'_i=x_i+\xi_i(x)t+\frac{t^2}{1\cdot2}\sum\limits_k^{1,\ldots,n} \xi_k\frac{\partial\xi_i}{\partial x_k}+\cdots .
\]
We can write this in a more expressive way by introducing, for any function $f(x_1,x_2,\ldots,x_n)$ of the variables $x$, the operation symbol
\[
 Xf=\sum\limits_i^{1,\ldots,n}\xi_i\frac{\partial f}{\partial x_i}
\]
and by denoting by $X^2f$, $X^3f,\ldots$ the result of the same operation repeated two, three, \ldots times; the preceding expansion then takes the form
\begin{equation*}\tag{17}
 x'_i=x_i+tXx_i+\frac{t^2}{1\cdot2}X^2x_i+\frac{t^3}{1\cdot2\cdot3}X^3x_i+\cdots .
\end{equation*}
More generally, if we wish to expand in powers of $t$ any function $f(x'_1,x'_2,\ldots,x'_n)$, regular at the point $(x_1,x_2,\ldots,x_n)$, we shall have
\begin{equation*}\tag{18}
\begin{aligned}
 f(x'_1,x'_2,\ldots,x'_n)&=f(x_1,x_2,\ldots,x_n)+tXf+\frac{t^2}{1\cdot2}X^2f+\cdots\\
 &\quad+\frac{t^s}{1\cdot2\cdots s}X^sf+\cdots .
\end{aligned}
\end{equation*}
Sometimes, for brevity, we shall denote this expansion by the symbol
\[
 f(x'_1,x'_2,\ldots,x'_n)=e^{tX}f,
\]
where $e^{tX}$ is to be formally expanded by the exponential series, and the powers of $X$ are here to be considered as operating on $f$.

Consider now that transformation of the group $G_1$ which corresponds to an infinitely small value $\delta t$ of the parameter $t$. This infinitesimal transformation of the group, when infinitesimals of order higher than $\delta t$ are neglected, gives the $x$ the infinitesimal increments
\[
 \delta x_i=\xi_i\,\delta t\qquad (i=1,2,\ldots,n),
\]
and therefore gives any function $f$ the increment
\[
 \delta f=Xf\,\delta t.
\]
It is clear that the effective expression of the symbol $Xf$ determines by itself completely the infinitesimal transformation, since it supplies the values of the $\xi$. Thus the symbol $Xf$ itself is regarded by Lie as the symbol of the transformation itself; that is, it denotes the infinitesimal transformation $Xf$. Naturally, if one changes the parameter $t$ into a new $\tau$ by putting
\[
 t=\varphi(\tau),
\]
the values of the $\xi$ are simultaneously multiplied by a constant factor, namely by the value of $\dfrac{dt}{d\tau}$ for $\tau=0$. Thus an infinitesimal transformation is not changed by multiplying it by any constant $c$. Observe now that the differential equations
\[
 \frac{dx'_i}{dt}=\xi_i(x'_1,x'_2,\ldots,x'_n),
\]
interpreted geometrically in the space of $n$ dimensions $(x'_1,x'_2,\ldots,x'_n)$, determine a system $\infty^{n-1}$ of regular curves; these are called the trajectories of the group, or also, according to Monge, the characteristics of the partial differential equation
\[
 Xf=0.
\]
We may therefore say: the symbol $Xf$ of the infinitesimal transformation is nothing other than the derivative of the function $f$ calculated in the direction of the trajectory of the group.

Having thus introduced the concept of infinitesimal transformation, we see that every group $G_1$ contains an infinitesimal transformation $Xf$, determined up to a constant factor; this, in turn, completely determines the group, which is generated from $Xf$ according to the formulae (17) or (18). The group itself is therefore indicated simply as the group $Xf$.

We state again the theorem at the end of \S{}9 in the new form:
Every group with one parameter
\[
 x'_i=f_i(x_1,x_2,\ldots,x_n;a)
\]
determines an infinitesimal transformation $Xf$, and all its transformations are obtained by composing one fixed transformation of it with the transformations of the group $G_1$, containing the identity, generated by $Xf$.

\section*{\S{} 11. Change of variables in infinitesimal transformations}

It has been seen how a group $G_1$, containing the identity, is perfectly determined by its infinitesimal transformation $Xf$. Now if, by means of a change of independent variables expressed by the formulae
\begin{equation*}\tag{19}
 y_i=\varphi_i(x_1,x_2,\ldots,x_n),
\end{equation*}
which pass from the old variables $x$ to the new variables $y$, the group $G_1$ is transformed into a similar group (\S{}9) $\Gamma_1$, one may ask how the infinitesimal transformation $Yf$ of the group $\Gamma_1$ is to be calculated. It is easy to see that the new symbol $Yf$ will be obtained simply by expressing, in the old $Xf$, the $x$ in terms of the $y$. For if we suppose
\[
 f(x_1,x_2,\ldots,x_n)=F(y_1,y_2,\ldots,y_n),
\]
we shall have
\[
 Xf=\sum\limits_\nu^{1,\ldots,n}\frac{\partial F}{\partial y_\nu}\,Xy_\nu.
\]
Having expressed the $Xy_\nu$ in terms of the $y$, put
\[
 Xy_\nu=\eta_\nu(y_1,y_2,\ldots,y_n),
\]
and
\[
 YF=\sum\limits_\nu^{1,\ldots,n}\eta_\nu\frac{\partial F}{\partial y_\nu};
\]
whatever $f$ may be, one has identically
\[
 Xf=YF,
\]
provided the $x,y$ are linked by the transformation formulae (19). Therefore one also has
\[
 X^2f=Y^2F,
\]
and, in general,
\[
 X^sf=Y^sF.
\]
Consequently (18) becomes
\[
\begin{aligned}
 F(y'_1,y'_2,\ldots,y'_n)&=F(y_1,y_2,\ldots,y_n)+tYF+\frac{t^2}{1\cdot2}Y^2F+\cdots\\
 &=e^{tY}F,
\end{aligned}
\]
a formula which holds for any $F$ and proves the stated property.

Thus every finite transformation
\[
 x'_i=x_i+tXx_i+\frac{t^2}{1\cdot2}X^2x_i+\cdots
\]
of $G_1$ is changed, for the transformed group, into the other
\[
 y'_i=y_i+tYy_i+\frac{t^2}{1\cdot2}Y^2y_i+\cdots,
\]
the value of $t$ being the same in the two formulae.

\section*{\S{} 12. The series $\infty^r$ of transformations and the corresponding infinitesimal transformations}

The object of the investigations we are now going to make is to extend to finite continuous groups with several parameters the results obtained above for groups with one parameter. In particular we shall prove that a group with $r$ parameters, which contains the identity and together with every transformation its inverse, is generated by $r$ infinitesimal transformations
\[
 X_1f,\ X_2f,\ldots,\ X_rf
\]
which are linearly independent of one another, in the sense that no identity of the form
\[
 e_1X_1f+e_2X_2f+\cdots+e_rX_rf=0,
\]
can occur, where the $e$ are numerical constants.

To establish these fundamental results it is useful to prefix some general considerations on series $\infty^r$ of transformations in general:
\begin{equation*}\tag{20}
 x'_i=f_i(x_1,x_2,\ldots,x_n; a_1,a_2,\ldots,a_r)
\end{equation*}
and on the infinitesimal transformations which they determine. Thus, among other things, we shall again establish the differential equations (A) of \S{}7, recognizing their intrinsic meaning.

Take any transformation $S_a$\footnote{As usual, we symbolically denote by $S_a$ that transformation which corresponds to the values $a_1,a_2,\ldots,a_r$ of the parameters in (20).} of the series (20), and an infinitely near transformation $S_{a+\delta a}$ of the same series, where the parameters $a$ have received infinitesimal increments $\delta a$. If we observe that
\[
 S_{a+\delta a}=S_a(S_a^{-1}S_{a+\delta a}),
\]
\[
 S_{a+\delta a}=(S_{a+\delta a}S_a^{-1})S_a,
\]
then we see that one can pass from $S_a$ to the infinitely near $S_{a+\delta a}$ either by first carrying out $S_a$ and then the infinitesimal transformation $S_a^{-1}S_{a+\delta a}$, or also by first carrying out the infinitesimal transformation $S_{a+\delta a}S_a^{-1}$ and then $S_a$. With respect to these two infinitesimal transformations, observe that in general they do not belong to the series (20), and that one is the transform of the other by means of $S_a$, as is seen from the formula
\[
 S_a^{-1}S_{a+\delta a}=S_a^{-1}(S_{a+\delta a}S_a^{-1})S_a.
\]

The effective analytic expressions of these two infinitesimal transformations can be written easily. Indeed, suppose that (20), solved with respect to the $x$, gives
\begin{equation*}\tag{20*}
 x_i=F_i(x'_1,x'_2,\ldots,x'_n; a_1,a_2,\ldots,a_r).
\end{equation*}
The first transformation $S_a^{-1}S_{a+\delta a}$ will be given by
\[
\begin{aligned}
(S_a^{-1}S_{a+\delta a})\quad x'_i={}&f_i\bigl(F_1(x,a),F_2(x,a),\ldots,F_n(x,a);\\
& a_1+\delta a_1,a_2+\delta a_2,\ldots,a_r+\delta a_r\bigr),
\end{aligned}
\]
and the second by
\[
\begin{aligned}
(S_{a+\delta a}S_a^{-1})\quad x'_i={}&F_i\bigl(f_1(x,a+\delta a),f_2(x,a+\delta a),\ldots,f_n(x,a+\delta a);\\
& a_1,a_2,\ldots,a_r\bigr).
\end{aligned}
\]
Since (20) and (20*) are inverse to one another, one can write for this
\begin{equation*}\tag{21}
(S_a^{-1}S_{a+\delta a})\quad
 x'_i=x_i+\sum\limits_k^{1,\ldots,r}\left[\frac{\partial f_i(\xi,a)}{\partial a_k}\right]_{\xi=F(x,a)}\delta a_k,
\end{equation*}
and
\[
(S_{a+\delta a}S_a^{-1})\quad
 x'_i=x_i+\sum\limits_k^{1,\ldots,r}\delta a_k\sum\limits_\nu^{1,\ldots,n}\frac{\partial f_\nu(x,a)}{\partial a_k}
 \left[\frac{\partial F_i(\xi,a)}{\partial \xi_\nu}\right]_{\xi=f(x,a)}.
\]
The expression of the second can be simplified by observing that, from the identity
\[
 F_i(f_1(x,a),f_2(x,a),\ldots,f_n(x,a))=x_i,
\]
differentiating with respect to $a_k$, we deduce
\[
 \left[\frac{\partial F_i(\xi,a)}{\partial a_k}\right]_{\xi=f(x,a)}+
 \sum\limits_\nu^{1,\ldots,n}\left[\frac{\partial F_i(\xi,a)}{\partial \xi_\nu}\right]_{\xi=f(x,a)}\frac{\partial f_\nu(x,a)}{\partial a_k}=0,
\]
whence it follows, for $S_{a+\delta a}S_a^{-1}$, that
\begin{equation*}\tag{21*}
(S_{a+\delta a}S_a^{-1})\quad
 x'_i=x_i-\sum\limits_k^{1,\ldots,r}\left[\frac{\partial F_i(\xi,a)}{\partial a_k}\right]_{\xi=f(x,a)}\delta a_k.
\end{equation*}
If we imagine the $a$ fixed and the $\delta a$ variable arbitrarily, (21), (21*) coordinate with $S_a$ two series of infinitesimal transformations by means of which one passes from $S_a$ to any other infinitely near transformation in the series (20). Now observe that (21) is written, with the notation of infinitesimal transformations, as
\[
 Z\Phi=\sum\limits_k^{1,\ldots,r}\delta a_k\sum\limits_i^{1,\ldots,n}\left[\frac{\partial f_i(\xi,a)}{\partial a_k}\right]_{\xi=F(x,a)}\frac{\partial\Phi}{\partial x_i},
\]
and is therefore a linear combination of the $r$ infinitesimal transformations
\begin{equation*}\tag{22}
 Z_k\Phi=\sum\limits_i^{1,\ldots,n}\left[\frac{\partial f_i(\xi,a)}{\partial a_k}\right]_{\xi=F(x,a)}\frac{\partial\Phi}{\partial x_i}.
\end{equation*}
These, for arbitrary values of the $a$, are certainly linearly independent of one another; otherwise the $r$ parameters $a$ would not all be essential. The same naturally holds for the series (21*), which is nothing other than a transform of the series (21). Conversely, it is clear that in any series (21) of transformations, if the $r$ infinitesimal transformations $Z_k$ given by (22) are independent, the $r$ parameters $a$ will be essential.

Finally observe that the generic infinitesimal transformation (21)
\[
 Z\Phi=\sum\limits_k^{1,\ldots,r}\delta a_k Z_k\Phi,
\]
depending on the $r-1$ ratios of the $\delta a$, runs through a series $\infty^{r-1}$ of distinct infinitesimal transformations.

Summing up the results obtained, we therefore have:
Every series (20), $\infty^r$ of transformations, coordinates with each of its transformations $S_a$ two series (21), (21*) of $\infty^{r-1}$ infinitesimal transformations, by means of which one passes, in the indicated way, from $S_a$ to all the infinitely near transformations in the same series.


\section*{\S{} 13. Particular case of a group}

Suppose now that the series $\infty^r$ given by (20) constitutes a group. Then any two transformations $S_a,S_b$ of the series compose into a third transformation $S_c$ of the same series:
\begin{equation*}\tag{23}
 S_aS_b=S_c,
\end{equation*}
where the parameters $c$ are certain functions of $a,b$,
\begin{equation*}\tag{24}
 c_d=\chi_d(a_1,a_2,\ldots,a_r; b_1,b_2,\ldots,b_r),\qquad d=1,2,\ldots,r,
\end{equation*}
and these relations, as we showed in \S{}6, are solvable both with respect to the $a$ and with respect to the $b$. This also follows at once from the symbolic equality (23): if we keep the $a$ fixed and let the $b$ vary arbitrarily, $S_c$ will run, like $S_b$, through a series $\infty^r$ of transformations, and therefore the $r$ functions $c$ are independent with respect to the $b$, that is,
\[
 \frac{\partial(c_1,c_2,\ldots,c_r)}{\partial(b_1,b_2,\ldots,b_r)}\ne0;
\]
and similarly for
\[
 \frac{\partial(c_1,c_2,\ldots,c_r)}{\partial(a_1,a_2,\ldots,a_r)}.
\]
Now give arbitrary increments $\delta a,\delta b$ to both the $a$ and the $b$. It will suffice to calculate the corresponding increments $\delta c$ from (24), by the formulae
\[
 \delta c_d=\sum\limits_k^{1,\ldots,r}\left(\frac{\partial\chi_d}{\partial a_k}\delta a_k+\frac{\partial\chi_d}{\partial b_k}\delta b_k\right),
\]
and one will have
\[
 S_{a+\delta a}S_{b+\delta b}=S_{c+\delta c}.
\]
Let us now take the $\delta a,\delta b$ so related to one another that the $\delta c$ all vanish, that is,
\begin{equation*}\tag{25}
 \sum\limits_k^{1,\ldots,r}\frac{\partial\chi_\lambda}{\partial a_k}\delta a_k+
 \sum\limits_k^{1,\ldots,r}\frac{\partial\chi_\lambda}{\partial b_k}\delta b_k=0,
 \qquad \lambda=1,2,\ldots,r.
\end{equation*}
Then it will follow that
\[
 S_{a+\delta a}S_{b+\delta b}=S_c=S_aS_b,
\]
that is,
\[
 S_a^{-1}S_{a+\delta a}\, S_{b+\delta b}S_b^{-1}=1.
\]
This proves that $S_{b+\delta b}S_b^{-1}$ is the inverse of $S_a^{-1}S_{a+\delta a}$. But clearly the inverse of $S_{b+\delta b}S_b^{-1}$ is obtained by changing the signs of all the $\delta b$; hence, as a consequence of (24), we have
\begin{equation*}\tag{26}
 S_a^{-1}S_{a+\delta a}=S_{b-\delta b}S_b^{-1}.
\end{equation*}
Observe now that in (25) we can choose the $\delta a$ arbitrarily and derive the $\delta b$, or conversely. It follows that the transformation on the left in (26) can be identified with an arbitrary one of the series (21), and likewise that on the right with an arbitrary one of the series (21*). Thus in this case the two series of transformations coincide. Moreover, if we keep the $b$ fixed and let the $a$ vary arbitrarily, (26) shows that the first series of infinitesimal transformations (21), coordinated with $S_a$, remains the same as $S_a$ varies, and the same naturally holds for the second series (21*).

We therefore conclude: with every group of transformations in $r$ parameters there is coordinated a perfectly determined series of $\infty^{r-1}$ infinitesimal transformations, by which, multiplying $S_a$ on the right as well as on the left, one obtains all transformations of the group infinitely near to $S_a$.

\section*{\S{} 14. New derivation of the fundamental differential equations (A) of \S{}7}

Let us now translate the symbolic equality (26) into formulae, substituting for the infinitesimal transformations of the two sides their effective expressions given by (21), (21*). We shall have:
\begin{equation*}\tag{27}
 \sum\limits_k^{1,\ldots,r}\left[\frac{\partial f_i(\xi,a)}{\partial a_k}\right]_{\xi=F(x,a)}\delta a_k
 =
 \sum\limits_j^{1,\ldots,r}\left[\frac{\partial F_i(\xi,b)}{\partial b_j}\right]_{\xi=f(x,b)}\delta b_j,
 \qquad i=1,2,\ldots,n,
\end{equation*}
relations that are a consequence of (25). Now, solving the equations (25) with respect to the $\delta b$, we write
\begin{equation*}\tag{28}
 \delta b_j=\sum\limits_k^{1,\ldots,r}\psi_{jk}(a,b)\delta a_k,
\end{equation*}
and the $\psi_{jk}$ will have precisely the meaning assigned to them in \S{}7, since they represent the derivatives of the $b$ with respect to the $a$ when the $c$ are kept fixed. Substituting the values (28) in (27), and remembering that the $\delta a$ are independent, we shall deduce
\begin{equation*}\tag{29}
 \left[\frac{\partial f_i(\xi,a)}{\partial a_k}\right]_{\xi=F(x,a)}
 =
 \sum\limits_j^{1,\ldots,r}\left[\frac{\partial F_i(\xi,b)}{\partial b_j}\right]_{\xi=f(x,b)}\psi_{jk}(a,b).
\end{equation*}
If we now introduce the notation
\begin{equation*}\tag{a}
 \left[\frac{\partial F_i(\xi,b)}{\partial b_j}\right]_{\xi=f(x,b)}=\Phi_{ji}(x,b),
\end{equation*}
and then in (29) change $x_i$ into $x'_i=f_i(x,a)$, we obtain the following equations:
\begin{equation*}\tag{30}
 \frac{\partial x'_i}{\partial a_k}=\sum\limits_j^{1,\ldots,r}\Phi_{ji}(x',b)\psi_{jk}(a,b).
\end{equation*}
These are the very equations (11) that we have already established in \S{}7, since not only do the $\psi_{jk}(a,b)$ have the meaning assigned to them in \S{}7, but also the $\Phi_{ji}(x',b)$, as follows from the following observations. One has identically
\begin{equation*}\tag{b}
 \left[F_i(\xi,b)\right]_{\xi=f(x,b)}=x_i,
\end{equation*}
and therefore, differentiating with respect to $b_j$,
\[
 \left[\frac{\partial F_i(\xi,b)}{\partial b_j}\right]_{\xi=f(x,b)}
 =-\sum\limits_\lambda^{1,\ldots,n}\left[\frac{\partial F_i(\xi,b)}{\partial\xi_\lambda}\right]_{\xi=f(x,b)}
 \frac{\partial f_\lambda(x,b)}{\partial b_j},
\]
whence, by (a),
\[
 \Phi_{ji}(x,b)= -\sum\limits_\lambda^{1,\ldots,n}\left[\frac{\partial F_i(\xi,b)}{\partial\xi_\lambda}\right]_{\xi=f(x,b)}
 \frac{\partial f_\lambda(x,b)}{\partial b_j}.
\]
It will therefore suffice, in order to prove the asserted coincidence, to show that
\[
 \left[\frac{\partial F_i(\xi,b)}{\partial\xi_\lambda}\right]_{\xi=f(x,b)}=\Theta_{\lambda i},
\]
where $\Theta_{\lambda i}$ is the quotient of the algebraic complement of $\dfrac{\partial f_\lambda}{\partial x_i}$ in the functional determinant
\[
 J=\frac{\partial(f_1,f_2,\ldots,f_n)}{\partial(x_1,x_2,\ldots,x_n)},
\]
divided by $J$ itself. But this follows at once from (b), differentiated with respect to any $x$, for example $x_\ell$, which gives\footnote{Recall the meaning of $\varepsilon_{\ell i}$: $\varepsilon_{\ell i}=0$ for $\ell\ne i$, and $\varepsilon_{\ell i}=1$ for $\ell=i$.}
\[
 \sum\limits_\lambda^{1,\ldots,n}\left[\frac{\partial F_i(\xi,b)}{\partial\xi_\lambda}\right]_{\xi=f(x,b)}
 \frac{\partial f_\lambda(x,b)}{\partial x_\ell}=\varepsilon_{\ell i}.
\]
Having thus obtained again (30), which coincide with (11) of \S{}7, it will suffice once more to give the $b$ arbitrary fixed values in order to deduce the fundamental differential equations (A) of \S{}7:
\begin{equation*}\tag{A}
 \frac{\partial x'_i}{\partial a_k}=\sum\limits_s^{1,\ldots,r}\xi_{si}(x')\psi_{sk}(a),
 \qquad
 \begin{cases}
 i=1,2,\ldots,n,\\
 k=1,2,\ldots,r.
 \end{cases}
\end{equation*}
With the new derivation these appear under a new aspect, as the translation into formulae of the symbolic equation (26), to which they are equivalent.

\section*{\S{} 15. The $r$ infinitesimal transformations $X_\lambda f=\sum\limits_i^{1,\ldots,n} \xi_{\lambda i}\dfrac{\partial f}{\partial x_i}$ $(\lambda=1,2,\ldots,r)$}

If we merely suppose that the series $\infty^r$ of transformations (20) verifies the system of differential equations (A), as it certainly will in particular if the series constitutes a group, we shall be able to put in another form the $r$ infinitesimal transformations of \S{}12
\[
 Z_1\Phi,
 Z_2\Phi,
 \ldots,
 Z_r\Phi,
\]
by which all the infinitesimal transformations (21) that carry a transformation $S_a$ to the infinitely near $S_{a+\delta a}$ are linearly combined.

Since (A) are satisfied, the increments $\delta x_i$ which the $x$ receive under an arbitrary infinitesimal transformation (21) are given by
\[
\begin{aligned}
 \delta x_i
 &=\sum\limits_k^{1,\ldots,r}\left[\frac{\partial f_i(\xi,a)}{\partial a_k}\right]_{\xi=F(x,a)}\delta a_k\\
 &=\sum\limits_k^{1,\ldots,r}\delta a_k\sum\limits_s^{1,\ldots,r}\xi_{si}(x)\psi_{sk}(a).
\end{aligned}
\]
The most general infinitesimal transformation under discussion is therefore written
\[
 Xf=\sum\limits_{s,k}^{1,\ldots,r}\delta a_k\psi_{sk}(a)\sum\limits_i^{1,\ldots,n}\xi_{si}\frac{\partial f}{\partial x_i},
\]
or, putting
\[
 \sum\limits_k^{1,\ldots,r}\psi_{sk}(a)\delta a_k=\lambda_s,
\]
\[
 Xf=\sum\limits_s^{1,\ldots,r}\lambda_s\sum\limits_i^{1,\ldots,n}\xi_{si}\frac{\partial f}{\partial x_i},
\]
where the $\lambda$ are arbitrary constants. If we put
\begin{equation*}\tag{31}
 X_1f=\sum\limits_i^{1,\ldots,n}\xi_{1i}\frac{\partial f}{\partial x_i},\quad
 X_2f=\sum\limits_i^{1,\ldots,n}\xi_{2i}\frac{\partial f}{\partial x_i},\quad \ldots,\quad
 X_rf=\sum\limits_i^{1,\ldots,n}\xi_{ri}\frac{\partial f}{\partial x_i},
\end{equation*}
we see that $Xf$ is a linear combination, with constant coefficients, of these $r$ fundamental infinitesimal transformations
\[
 X_1f,
 X_2f,
 \ldots,
 X_rf.
\]
These are linearly independent, as follows from what we have already said in \S{}12, or also from the observations at the end of \S{}8, from which it follows that no identity of the form
\[
 e_1X_1f+e_2X_2f+\cdots+e_rX_rf=0
\]
with the $e$ constant can subsist.

Thus: If a series $\infty^r$ of transformations satisfies the differential equations (A), one passes from an arbitrary transformation $S_a$ of the series to an infinitely near one $S_{a+\delta a}$ by composing $S_a$ with an infinitesimal transformation $\sum\limits_s^{1,\ldots,r}\lambda_sX_sf$, a linear combination with constant coefficients $\lambda$ of the $r$ fixed, linearly independent infinitesimal transformations
\[
 X_sf=\sum\limits_i^{1,\ldots,n}\xi_{si}\frac{\partial f}{\partial x_i}\qquad(s=1,2,\ldots,r).
\]
If we recall (\S{}10) that an infinitesimal transformation generates a group with one parameter, it is natural to ask now whether the property just proved for the infinitesimal transformation $\sum\limits_s^{1,\ldots,r}\lambda_sX_sf$ also holds for the finite transformations of the group $G_1$ with one parameter generated by it; that is, whether by composing $S_a$ with a finite transformation of $G_1$ one still obtains a transformation of the primitive series.

The answer is affirmative, as we now proceed to show; and from this result consequences of fundamental importance for the theory of groups will be deduced.

\section*{\S{} 16. Fundamental theorem for the series of transformations that satisfy (A)}

To prove what we have asserted above, we begin by observing that, since for generic values of the $a$ the determinant of the $\psi$ is different from zero, we can, as we already did in \S{}8, solve (A) with respect to the $\xi$, and write them in the equivalent form
\begin{equation*}\tag{B}
 \xi_{ji}(x')=\sum\limits_k^{1,\ldots,r}\alpha_{jk}(a)\frac{\partial x'_i}{\partial a_k}.
\end{equation*}
The infinitesimal transformation $\sum\limits_s^{1,\ldots,r}\lambda_sX_sf$, written out, is given by
\[
 \delta x_i=\delta t\sum\limits_s^{1,\ldots,r}\lambda_s\xi_{si}(x),\qquad (i=1,2,\ldots,n),
\]
and, calculated from (A), is written
\[
 \delta x_i=\sum\limits_{s,k}^{1,\ldots,r}\xi_{si}(x)\psi_{sk}(a)\delta a_k,
\]
whence, comparing, one has
\[
 \sum\limits_s^{1,\ldots,r}\xi_{si}\left(\lambda_s\delta t-\sum\limits_k^{1,\ldots,r}\psi_{sk}(a)\delta a_k\right)=0
 \qquad(i=1,2,\ldots,n),
\]
and therefore\footnote{Otherwise $X_1f,X_2f,\ldots,X_rf$ would not be independent.}
\[
 \sum\limits_k^{1,\ldots,r}\psi_{sk}(a)\delta a_k=\lambda_s\delta t.
\]
Solving with respect to the $\delta a_k$, we have the formulae
\begin{equation*}\tag{c}
 \delta a_k=\delta t\sum\limits_s^{1,\ldots,r}\alpha_{sk}(a)\lambda_s,
\end{equation*}
by which the increments of the $a$ are to be calculated in the passage from $S_a$ to $S_{a+\delta a}$ which results from composing $S_a$ with the infinitesimal transformation $\sum\limits_s^{1,\ldots,r}\lambda_sX_sf$.

Fix for the $a$ a system of particular values
\[
 a_1^{(0)},a_2^{(0)},\ldots,a_r^{(0)}
\]
subject only to the condition that for them the determinant of the $\psi_{jk}(a^{(0)})$, and hence that of the $\alpha_{jk}(a^{(0)})$, is not zero. Having then fixed arbitrary constant values for the $\lambda$, determine the $a$ as functions of $t$ according to (c), from the differential equations
\begin{equation*}\tag{32}
 \frac{da_k}{dt}=\sum\limits_s^{1,\ldots,r}\lambda_s\alpha_{sk}(a)
\end{equation*}
with initial conditions
\[
 a_k=a_k^{(0)}\qquad\text{for }t=0.
\]
The transformation $S_a$ corresponding to the variation of $t$ will run through as many transformations of the series (20) in a neighborhood of $S_{a^{(0)}}$. It is indeed easy to see that, by then varying the $\lambda$, $S_a$ will run through all the transformations in the neighborhood of $S_{a^{(0)}}$. For, from the form of (32), it is evident that the integrals $a$ will depend on the variables
\[
 u_1=\lambda_1t,
 \quad u_2=\lambda_2t,
 \ldots,
 \quad u_r=\lambda_rt;
\]
and since the functional determinant of the $a$ with respect to the $u$ reduces, for $t=0$, precisely to $\lvert\alpha_{ik}(a^{(0)})\rvert\ne0$, it follows that we can conversely express $u_1,u_2,\ldots,u_r$ in terms of $a_1,a_2,\ldots,a_r$, so that by suitably choosing the $\lambda$ and $t$ the parameters $a$ will receive arbitrary values.

Now if, as above, we keep the $\lambda$ fixed and let $t$ vary, $S_a$ will run through a series $\infty^1$ of transformations within the complete series
\[
 x'_i=f_i(x,a).
\]
The $x'_i$, considered as functions of $t$, will, because of the fundamental equations (A) and (32), satisfy the differential equations
\[
 \frac{dx'_i}{dt}=\sum\limits_s^{1,\ldots,r}\xi_{si}(x')\sum\limits_k^{1,\ldots,r}\psi_{sk}(a)\sum\limits_j^{1,\ldots,r}\lambda_j\alpha_{jk}(a),
\]
or
\[
 \frac{dx'_i}{dt}=\sum\limits_s^{1,\ldots,r}\lambda_s\xi_{si}(x').
\]
According to the result obtained for groups with one parameter at the end of \S{}9, the transformation $S_a$ is therefore produced by first performing the fixed transformation $S_{a^{(0)}}$ and then a finite transformation of the group $G_1$ containing the identity, generated by the infinitesimal transformation $\sum\limits_j^{1,\ldots,r}\lambda_jX_jf$; and this is exactly what we wished to prove.

We can therefore state the theorem.
If a series $\infty^r$ of transformations satisfies a system of differential equations of type (A), and if $S_{a^{(0)}}$ is a particular transformation of it for which the determinant of the $\psi_{hk}(a^{(0)})$ is not zero, then every transformation of the series in a neighborhood of $S_{a^{(0)}}$ is obtained by composing $S_{a^{(0)}}$ with a finite transformation of the group generated by an arbitrary linear combination, with constant coefficients, of the $r$ independent infinitesimal transformations
\[
 X_\lambda f=\sum\limits_i^{1,\ldots,n}\xi_{\lambda i}\frac{\partial f}{\partial x_i}.
\]
Conversely, however the $\lambda$ are taken, the composition of $S_{a^{(0)}}$ with a finite transformation of the group $\sum\limits_j^{1,\ldots,r}\lambda_jX_jf$ gives rise to a transformation of the same series.

\chapter*{Chapter II}
\addcontentsline{toc}{chapter}{Chapter II}

\begin{center}
\begin{minipage}{0.88\linewidth}
\centering
The three fundamental theorems of the theory of groups.\par
Mixed groups.
\end{minipage}
\end{center}

\section*{\S{} 17. The first fundamental theorem}

The whole theory of continuous groups was founded by Lie on the following three successive theorems, with which we shall now concern ourselves.

The first of them, to which we have already alluded elsewhere, gives in a certain sense the inverse of the final theorem of \S{}7, by establishing that the differential equations of type (A), together with certain secondary conditions, are characteristic for the series that constitute a group.

For series with one parameter we have already established the theorem in question, and now we wish to extend it to the case of a series with $r$ parameters
\begin{equation*}\tag{1}
 x'_i=f_i(x_1,x_2,\ldots,x_n; a_1,a_2,\ldots,a_r),\qquad i=1,2,\ldots,n,
\end{equation*}
which we suppose to satisfy a system of differential equations of type (A)
\begin{equation*}\tag{A}
 \frac{\partial x'_i}{\partial a_k}=\sum\limits_s^{1,\ldots,r}\xi_{si}(x')\psi_{sk}(a),
 \qquad
 \begin{cases}
 i=1,2,\ldots,n,\\
 k=1,2,\ldots,r.
 \end{cases}
\end{equation*}
For this it is enough to apply the results obtained in the preceding section. Let us begin by supposing that the series (1) contains the identity transformation, corresponding to the values $a_i=a_i^{(0)}$ of the parameters, and also that for these values of the parameters, which give $S_{a^{(0)}}=1$, the determinant of the $\psi$ is different from zero. Applying the theorem at the end of \S{}16, taking $S_a$ to be $S_{a^{(0)}}=1$, we see first that the series (1), at least in a certain neighborhood of the identity, is nothing but the set of all transformations of the $\infty^{r-1}$ groups $G_1$ generated by the infinitesimal transformations
\[
 \sum\limits_i^{1,\ldots,r} \lambda_i X_i f.
\]
If we then apply the theorem once more, taking for $S_a$ a transformation in the neighborhood of $S_{a^{(0)}}=1$, and denote by $S_e$ a transformation of one of the groups $G_1$, which is also an arbitrary transformation of the series (1), we obtain
\[
 S_aS_e=S_c,
\]
where $S_c$ is a new transformation of the series. Thus it is precisely shown that, under the assumed conditions, the series (1) constitutes a group and indeed a group with pairs of inverse transformations, since this property belongs, by \S{}9, to the groups $G_1$ that compose our group with $r$ parameters.

Summarizing the results now obtained with those established directly in \S{}7, we can give the first fundamental theorem the following statement.

If a series $\infty^r$ of transformations
\[
 x'_i=f_i(x_1,x_2,\ldots,x_n; a_1,a_2,\ldots,a_r),\qquad i=1,2,\ldots,n,
\]
forms a group, then the $x'_i$, considered as functions of the $a$, satisfy a system of differential equations of the form
\begin{equation*}\tag{A}
 \frac{\partial x'_i}{\partial a_k}=\sum\limits_s^{1,\ldots,r}\xi_{si}(x')\psi_{sk}(a),
 \qquad
 \begin{cases}
 i=1,2,\ldots,n,\\
 k=1,2,\ldots,r,
 \end{cases}
\end{equation*}
where the determinant of the $\psi_{sk}(a)$ is not identically zero and the $r$ infinitesimal transformations
\[
 X_kf=\sum\limits_i^{1,\ldots,n} \xi_{ki}(x)\frac{\partial f}{\partial x_i}\qquad (k=1,2,\ldots,r)
\]
are linearly independent.

Conversely, if a series $\infty^r$ of transformations satisfies a system of equations of type (A), contains the identity transformation as the one corresponding to the parameter values $a_i=a_i^{(0)}$, and, moreover, the determinant $|\psi_{hk}(a^{(0)})|$ is nonzero, then the series itself is a group with pairs of inverse transformations and coincides with the totality of the transformations contained in the one-parameter groups $G_1$ generated by the $\infty^{r-1}$ linear combinations
\[
 \sum\limits_i^{1,\ldots,r} \lambda_i X_i f.
\]
Under the hypothesis of the last part of the theorem the group (with pairs of inverse transformations) will be denoted by $G_r$ and will be said to be generated by the $r$ infinitesimal transformations
\[
 X_1f,\ X_2f,\ldots,\ X_rf.
\]

\section*{\S{} 18. Further cases}

To be able to apply the preceding theorem in its second part, the determinant
\[
 |\psi_{hk}(a^{(0)})|
\]
for the values $a^{(0)}$ of the parameters that give the identity must be different from zero. But now let us suppose only that the series $\infty^r$ of transformations (1) satisfies a system (A), and let $S_{\bar a}$ be a particular transformation of the series such that
\[
 |\psi_{hk}(\bar a)|\ne0.
\]
The series
\[
 S_{\bar a}^{-1}S_a
\]
contains $\infty^r$ transformations and, for $a_i=\bar a_i$, gives the identity. On the other hand, since the explicit formulas for the series $S_{\bar a}^{-1}S_a$ are
\[
 x'_i=f_i\bigl(F_1(x,\bar a),F_2(x,\bar a),\ldots,F_n(x,\bar a); a_1,a_2,\ldots,a_r\bigr),
\]
these new functions $x'_i$ still satisfy the system (A), because they differ from the primitive functions only by replacing the $x_i$ with $F_i(x,\bar a)$. Hence the fundamental theorem applies to this series $S_{\bar a}^{-1}S_a$, and we have the result:

If a series $\infty^r$ of transformations (1) satisfies a system (A), and if $S_{\bar a}$ is a transformation such that the determinant $|\psi_{hk}(\bar a)|$ is nonzero, then the $\infty^r$ transformations $S_{\bar a}^{-1}S_a$, as the $a$ vary, constitute a group $G_r$ generated by the $r$ infinitesimal transformations
\[
 X_kf=\sum\limits_i^{1,\ldots,n} \xi_{ki}\frac{\partial f}{\partial x_i}\qquad (k=1,2,\ldots,r).
\]
Suppose now that the series $\infty^r$ of transformations (1) constitutes a group. It then follows by itself that these transformations satisfy a system of equations of type (A); and, having chosen in the group a fixed transformation $S_{\bar a}$ such that $|\psi_{hk}(\bar a)|\ne0$, the $\infty^r$ transformations $S_{\bar a}^{-1}S_a$ constitute a group $G_r$ with pairs of inverse transformations, generated by the infinitesimal transformations $X_1f,X_2f,\ldots,X_rf$. Now if, besides $S_{\bar a}$, the inverse $S_{\bar a}^{-1}$ also appears in the group, the said series $S_{\bar a}^{-1}S_a$ coincides with the group itself. Thus one has the important theorem.

Every group with $r$ parameters whose transformations are inverse in pairs\footnote{From what has been seen, it is enough that, together with a certain transformation $S_{\bar a}$, its inverse $S_{\bar a}^{-1}$ exist in the group, provided that the determinant $|\psi_{hk}(\bar a)|$ is different from zero; then the theorem in the text is applicable.} is a group $G_r$ generated by $r$ infinitesimal transformations
\[
 X_1f,\ X_2f,\ldots,\ X_rf,
\]
whose expressions are obtained in the familiar way by establishing the system of differential equations (A) satisfied by the transformations of the group.

The groups $G_r$ with pairs of inverse transformations are those that occur almost exclusively, and the simplicity of their theory consists principally in the fact that their study can be carried out on the $r$ infinitesimal transformations that generate them. It is only of these groups that we shall always speak in what follows, unless the contrary is explicitly stated.

\section*{\S{} 19. The complete system of partial differential equations equivalent to system (A)}

The groups on which all the interest of our studies will be concentrated are, as has been said, the groups $G_r$ generated by $r$ infinitesimal transformations
\[
 X_kf=\sum\limits_i^{1,\ldots,n}\xi_{ki}\frac{\partial f}{\partial x_i}\qquad (k=1,2,\ldots,r)
\]
that are linearly independent. Now these $r$ transformations
\[
 X_1f,\ X_2f,\ldots,\ X_rf
\]
cannot be chosen arbitrarily if they are to generate a group; they must satisfy certain special conditions, since otherwise the aggregate of one-parameter groups $G_1$
\[
 \lambda_1X_1f+\lambda_2X_2f+\cdots+\lambda_rX_rf
\]
would form only an $\infty^r$ series of transformations without the group character. The expression of these necessary and sufficient conditions for the generation of a group $G_r$ is the content of the second fundamental theorem, which Lie called the principal theorem (\emph{Hauptsatz}) as the most important of the three. We begin by finding these conditions as necessary. Assuming that the series (1) constitutes a group $G_r$, the $f_i(x,a)$, considered as functions of the $a$, must satisfy a system of total differential equations (A), which must be unlimitedly integrable (\S{}7). The integrability conditions of the system (A), which we now develop, give precisely the conditions sought.

We know that the integration of a system of total differential equations such as system (A) is perfectly equivalent to that of a complete system of homogeneous linear equations for a single unknown function of the independent variables and the primitive unknown functions. In our case the unknown function is
\[
 F(x'_1,x'_2,\ldots,x'_n; a_1,a_2,\ldots,a_r)
\]
of the $n+r$ variables $x',a$, which must satisfy the $r$ homogeneous linear equations
\begin{equation*}\tag{2}
 \frac{\partial F}{\partial a_k}+\sum\limits_{j=1}^r\sum\limits_{i=1}^n\xi_{ji}(x')\psi_{jk}(a)\frac{\partial F}{\partial x'_i}=0,
 \qquad k=1,2,\ldots,r.
\end{equation*}
For the system to be unlimitedly integrable, it is necessary and sufficient that the system (2) be complete. Now, for calculating the integrability conditions, it is more useful to write (2) in the following form. In the nonzero determinant of the $\psi$ we denote, as usual, by $\alpha_{jk}(a)$ the algebraic complement of $\psi_{jk}$ divided by $\psi$ itself. Multiplying (2) by $\alpha_{jk}$ and summing with respect to $k$ from $1$ to $r$, and giving $j$ its $r$ values, we replace (2) by the equivalent system
\begin{equation*}\tag{2*}
 \sum\limits_k^{1,\ldots,r}\alpha_{jk}(a)\frac{\partial F}{\partial a_k}+\sum\limits_i^{1,\ldots,n}\xi_{ji}(x')\frac{\partial F}{\partial x'_i}=0,
 \qquad j=1,2,\ldots,r.
\end{equation*}
Of the two sums that make up the first member, the second, whose coefficients contain only the $x'$, is nothing other than the expression $X_jf$, with the $x$ replaced by the $x'$, and we shall therefore denote it by
\[
 X'_jf=\sum\limits_i^{1,\ldots,n}\xi_{ji}(x')\frac{\partial F}{\partial x'_i}.
\]
The coefficients of the first sum also depend only on the variables $a$ with respect to which the derivatives of $F$ are taken. Introducing the symbol
\begin{equation*}\tag{3}
 A_jF=\sum\limits_k^{1,\ldots,r}\alpha_{jk}(a)\frac{\partial F}{\partial a_k},
\end{equation*}
the complete system (2) to be integrated, generated by the system of fundamental differential equations (A), is written
\begin{equation*}\tag{C'}
 A_jF+X'_jF=0,\qquad j=1,2,\ldots,r.
\end{equation*}

\section*{\S{} 20. Integrability conditions}

In order that the system (C') be complete, all the alternants
\[
 (A_j+X'_j,\ A_k+X'_k),\qquad j,k=1,2,\ldots,r,
\]
must be composed linearly with the first members of (C') themselves. Observing that the $x'$ do not enter the $A$ and the $a$ do not enter the $X'$, these conditions are written
\begin{equation*}\tag{4}
 (A_j,A_k)+(X'_j,X'_k)=\sum\limits_s^{1,\ldots,r} \gamma_{jks}(A_sF+X'_sF),
\end{equation*}
where the $\gamma_{jks}$ denote suitable functions of the $x'$ and of the $a$, which, as will now be seen, reduce to constants.

Meanwhile, since (4) must hold for an arbitrary function $F$ of the $x'$ and therefore also for an arbitrary function of the $a$ alone or of the $x'$ alone, we shall separately have
\begin{equation*}\tag{5}
 (A_j,A_k)=\sum\limits_s^{1,\ldots,r} \gamma_{jks} A_sF,
\end{equation*}
\begin{equation*}\tag{6}
 (X'_j,X'_k)=\sum\limits_s^{1,\ldots,r} \gamma_{jks} X'_sF.
\end{equation*}
Comparing in (5) the coefficients of $\partial F/\partial a_i$ on the two sides, we have
\[
 A_j(\alpha_{ki})-A_k(\alpha_{ji})=\sum\limits_s^{1,\ldots,r} \gamma_{jks}\alpha_{si},\qquad i=1,2,\ldots,r.
\]
These may be considered as $r$ linear equations in
\[
 \gamma_{jk1},\gamma_{jk2},\ldots,\gamma_{jkr},
\]
and since the determinant of the $\alpha$ is not zero, they can be solved with respect to the $\gamma_{jks}$ ($s=1,2,\ldots,r$). Hence it follows first that the $\gamma_{jks}$ do not contain the $x'$. Similarly, if in (6) we compare the coefficients of $\partial F/\partial x_i$ on the two sides, we have
\[
 X'_j(\xi_{ki})-X'_k(\xi_{ji})=\sum\limits_s^{1,\ldots,r}\gamma_{jks}\xi_{si}(x'),
\]
whence, differentiating with respect to an arbitrary $a$,
\[
 \sum\limits_s^{1,\ldots,r}\frac{\partial\gamma_{jks}}{\partial a_q}X'_sF=0;
\]
therefore the derivatives $\partial\gamma_{jks}/\partial a_q$ must all be zero, since
\[
 X'_1F,\ X'_2F,\ldots,\ X'_rF
\]
are linearly independent. It follows that the $\gamma_{jks}$, containing neither the $x'$ nor the $a$, are numerical constants; we denote them by $c_{jks}$.

Thus: The conditions of unlimited integrability of the fundamental equations (A) are written
\begin{equation*}\tag{$\alpha$}
 (X'_j,X'_k)=\sum\limits_s^{1,\ldots,r} c_{jks}X'_sF,
\end{equation*}
\begin{equation*}\tag{$\beta$}
 (A_j,A_k)=\sum\limits_s^{1,\ldots,r} c_{jks}A_sF,
\end{equation*}
where the $c_{jks}$ are numerical constants.

\section*{\S{} 21. Conditions $(\alpha),(\beta)$ as sufficient}

Suppose now conversely that the $r$ infinitesimal transformations, linearly independent in $x'_1,x'_2,\ldots,x'_n$,
\[
 X'_jF=\sum\limits_i^{1,\ldots,n}\xi_{ji}(x')\frac{\partial F}{\partial x'_i}
\]
satisfy $(\alpha)$, and that the $r$ infinitesimal transformations in the $a$,
\[
 A_jF=\sum\limits_k^{1,\ldots,r}\alpha_{jk}(a)\frac{\partial F}{\partial a_k},
\]
also linearly independent, satisfy $(\beta)$. Then the system (C') is complete, and the equivalent system (2), being solved with respect to the derivatives
\[
 \frac{\partial F}{\partial a_1},\ \frac{\partial F}{\partial a_2},\ldots,\frac{\partial F}{\partial a_r},
\]
is Jacobian in normal form. Fix a particular system of values of the $a$
\[
 a_1^{(0)},a_2^{(0)},\ldots,a_r^{(0)},
\]
such only that the determinant $|\psi_{ik}(a^{(0)})|$ is not zero, and consider the principal solutions
\[
 F_1,F_2,\ldots,F_n
\]
of the system (2), which for $a_i=a_i^{(0)}$ reduce respectively to
\[
 x'_1,x'_2,\ldots,x'_n.
\]
Solving the system of equations
\[
 F_i(x'_1,x'_2,\ldots,x'_n; a_1,a_2,\ldots,a_r)=x_i,\qquad i=1,2,\ldots,n,
\]
with respect to the $x'$ gives
\begin{equation*}\tag{1'}
 x'_i=f_i(x_1,x_2,\ldots,x_n; a_1,a_2,\ldots,a_r),
\end{equation*}
and these functions $f_i(x,a)$ satisfy, as we know, the system (A) of total differential equations. Since, moreover, the series $\infty^r$ of transformations contains the identity corresponding to the values $a^{(0)}$ of the $a$, and $|\psi_{ik}(a^{(0)})|\ne0$, it is enough to invoke the second part of the first fundamental theorem (\S{}17) to see that (1') defines a group $G_r$ generated by the $r$ infinitesimal transformations $X_1f,X_2f,\ldots,X_rf$.

\section*{\S{} 22. New proof of the first fundamental theorem}

Instead of resorting, as we have just done, to the second part of the first fundamental theorem, we may here derive from the considerations of the preceding section a new proof of this theorem due to Schur.

Suppose in fact: first, that the $\infty^r$ series of transformations
\begin{equation*}\tag{a}
 x'_i=f_i(x_1,x_2,\ldots,x_n; a_1,a_2,\ldots,a_r)
\end{equation*}
satisfies a system of differential equations of the form
\begin{equation*}\tag{b}
 \frac{\partial x'_i}{\partial a_k}=\sum\limits_s^{1,\ldots,r}\xi_{si}(x')\psi_{sk}(a);
\end{equation*}
second, that the identity is contained in (a), for $a_k=a_k^{(0)}$, and that the determinant $|\psi_{ik}(a^{(0)})|$ is not zero.

Then, since (b) forms an unlimitedly integrable system, the infinitesimal transformations
\[
 X_jf,\qquad A_jf
\]
will necessarily satisfy $(\alpha),(\beta)$. We now try to determine $r$ unknown functions $c_1,c_2,\ldots,c_r$ of $r$ variables $b_1,b_2,\ldots,b_r$ from the system of total differential equations
\begin{equation*}\tag{c}
 \frac{\partial c_\lambda}{\partial b_k}=\sum\limits_s^{1,\ldots,r}\alpha_{\lambda s}(c_1,c_2,\ldots,c_r)\psi_{sk}(b_1,b_2,\ldots,b_r),
 \qquad \lambda,k=1,2,\ldots,r.
\end{equation*}
This system (c) is unlimitedly integrable, because the corresponding homogeneous linear system in partial derivatives for an unknown function $\Theta(c_1,c_2,\ldots,c_r;b_1,b_2,\ldots,b_r)$ is written
\[
 -\frac{\partial\Theta}{\partial b_k}+\sum\limits_{\lambda,s}^{1,\ldots,r}\alpha_{\lambda s}(c)\psi_{sk}(b)\frac{\partial\Theta}{\partial c_\lambda}=0,
 \qquad k=1,2,\ldots,r;
\]
multiplying the preceding equation by $\alpha_{jk}(b)$ and summing from $k=1$ to $k=r$, we get the equivalent system
\[
 \sum\limits_k^{1,\ldots,r}\alpha_{jk}(b)\frac{\partial\Theta}{\partial b_k}+\sum\limits_\lambda^{1,\ldots,r}\alpha_{j\lambda}(c)\frac{\partial\Theta}{\partial c_\lambda}=0.
\]
Now this is a complete system because of $(\beta)$.

We therefore integrate (c) under the initial conditions
\[
 c_i=a_i\quad\text{for}\quad b_k=a_k^{(0)},
\]
and let
\begin{equation*}\tag{d}
 c_i=\theta_i(a_1,a_2,\ldots,a_r; b_1,b_2,\ldots,b_r)
\end{equation*}
be the corresponding integrals. If we put
\begin{equation*}\tag{e}
 x_i''=f_i(x_1,x_2,\ldots,x_n;c_1,c_2,\ldots,c_r),
\end{equation*}
then by (b) we shall have
\[
 \frac{\partial x_i''}{\partial b_k}
 =\sum\limits_\lambda^{1,\ldots,r} \frac{\partial x_i''}{\partial c_\lambda}\frac{\partial c_\lambda}{\partial b_k}
 =\sum\limits_{\lambda,s,t}^{1,\ldots,r}\xi_{si}(x'')\psi_{s\lambda}(c)\alpha_{\lambda t}(c)\psi_{tk}(b).
\]
But
\[
 \sum\limits_\lambda^{1,\ldots,r}\psi_{s\lambda}(c)\alpha_{\lambda t}(c)=\varepsilon_{st}=\begin{cases}0,&s\ne t,\\1,&s=t,\end{cases}
\]
and therefore
\[
 \frac{\partial x_i''}{\partial b_k}=\sum\limits_s^{1,\ldots,r}\xi_{si}(x'')\psi_{sk}(b).
\]
These are, for the $x_i''$ considered as functions of the $b$, precisely equations (b); and since for $b_k=a_k^{(0)}$ they reduce, by (e), to
\[
 f_i(x_1,x_2,\ldots,x_n; a_1,a_2,\ldots,a_r)=x_i',
\]
it follows by the uniqueness theorem that
\[
 x_i''=f_i(x_1',x_2',\ldots,x_n'; b_1,b_2,\ldots,b_r).
\]
The functional equations
\[
 f_i\bigl(f_1(x,a),f_2(x,a),\ldots,f_n(x,a); b_1,b_2,\ldots,b_r\bigr)
 =f_i(x_1,x_2,\ldots,x_n; c_1,c_2,\ldots,c_r)
\]
characteristic of a group are therefore satisfied.

One should observe that in this proof of the first fundamental theorem by Schur, the formulas that connect the parameters $c$ of the composed transformation $S_c=S_bS_a$ to those $a,b$ of the components are obtained by integrating the total differential equations (c), without its being necessary to know the finite equations (a) of the group. We shall later deal with the so-called parametric groups on the basis of this observation.

\section*{\S{} 23. The second fundamental theorem (principal theorem)}

In order that the infinitesimal transformations
\[
 X_1f,\ X_2f,\ldots,\ X_rf
\]
which are linearly independent should generate a group $G_r$, it is necessary and sufficient, as was seen at the end of \S{}21, that the conditions (2) and (3) be satisfied. But, since the group is already perfectly determined, that is to say exists, from the expressions of the $X$ alone, it is natural to ask whether the first conditions
\begin{equation*}\tag{2}
 (X_i,X_k)=\sum\limits_s^{1,\ldots,r} c_{iks}X_s f
\end{equation*}
are not already enough to assure its existence; that is, to infer that the $\infty^{r-1}$ groups $G_1$ generated by the linear combinations with constant coefficients of the $X$
\[
 \sum\limits_k^{1,\ldots,r} a_kX_kf
\]
constitute a group $G_r$. The answer to this question is affirmative, and to prove it it will suffice to establish, according to \S{}21, that if the $r$ linearly independent infinitesimal transformations
\begin{equation*}\tag{2*}
 X_kf=\sum\limits_i^{1,\ldots,n}\xi_{ki}(x)\frac{\partial f}{\partial x_i}
\end{equation*}
satisfy (2), one can always find another $r$ linearly independent infinitesimal transformations
\[
 A_if=\sum\limits_k^{1,\ldots,r}\alpha_{ik}(a_1,a_2,\ldots,a_r)
        \frac{\partial f}{\partial a_k},
 \qquad i=1,2,\ldots,r,
\]
in $r$ independent variables $a_1,a_2,\ldots,a_r$, such that the relations
\begin{equation*}\tag{3}
 (A_i,A_k)=\sum\limits_s^{1,\ldots,r}c_{iks}A_s f
\end{equation*}
also hold. To find the required $A_if$ one may proceed as follows. Write the expressions $X$ $r$ times in $r$ systems of independent variables
\[
 (x^{(1)}_1,x^{(1)}_2,\ldots,x^{(1)}_n),\quad
 (x^{(2)}_1,x^{(2)}_2,\ldots,x^{(2)}_n),\quad \ldots,
 (x^{(r)}_1,x^{(r)}_2,\ldots,x^{(r)}_n),
\]
putting
\[
 X^{(\lambda)}_i\Phi=
 \sum\limits_k^{1,\ldots,n}\xi_{ik}(x^{(\lambda)}_1,x^{(\lambda)}_2,\ldots,x^{(\lambda)}_n)
       \frac{\partial\Phi}{\partial x^{(\lambda)}_k},
 \qquad i,\lambda=1,2,\ldots,r,
\]
and put
\begin{equation*}\tag{8}
 W_i\Phi=X_i^{(1)}\Phi+X_i^{(2)}\Phi+\cdots+X_i^{(r)}\Phi,
\end{equation*}
where $\Phi$ is any function whatever of the $rn$ variables $x^{(\lambda)}_k$. First of all it is easy to prove that the system of $rn$ homogeneous linear equations in the partial derivatives of $\Phi$
\begin{equation*}\tag{9}
 W_1\Phi=0,
 \quad W_2\Phi=0,
 \quad \ldots,
 \quad W_r\Phi=0
\end{equation*}
is made up of equations all independent, that is, no identity of the form
\[
 \chi_1W_1\Phi+\chi_2W_2\Phi+\cdots+\chi_rW_r\Phi=0
\]
can hold, the $\chi$ being suitable functions of the $rn$ variables $x^{(\lambda)}$. Such an identity would indeed give, for every $\lambda$,
\begin{equation*}\tag{10}
 \chi_1\xi_{1k}(x^{(\lambda)})+
 \chi_2\xi_{2k}(x^{(\lambda)})+
 \cdots+
 \chi_r\xi_{rk}(x^{(\lambda)})=0,
 \qquad k=1,2,\ldots,n.
\end{equation*}
If we suppose them compatible, then in the first group of $n$ equations obtained by taking $\lambda=1$ there must be only $r-\nu$ independent equations. The following $n$ equations obtained by taking $\lambda=2$ cannot all be consequences of the preceding ones; otherwise we could satisfy the $n$ equations
\[
 \sum\limits_i^{1,\ldots,r}\chi_i\xi_{ik}(x^{(\lambda)}_1,x^{(\lambda)}_2,\ldots,x^{(\lambda)}_n)=0
\]
with values of the $\chi$ independent of the $x^{(2)}$, and hence the $r$ expressions
\[
 X_1^{(2)}\Phi,
 X_2^{(2)}\Phi,
 \ldots,
 X_r^{(2)}\Phi
\]
would not be independent. Thus the number of equations of the second group independent of one another and of the preceding ones is certainly $\nu_1$. In the same way one sees that each of the $r$ groups supplies at least one equation independent of the preceding ones; in this way one reaches $r$ independent equations, whence all the $\chi$ must vanish.

Now the system of $r$ equations (9) is a complete system, because, since in general
\[
 (X_i^{(\lambda)},X_k^{(\lambda)})=\sum\limits_s^{1,\ldots,r}c_{iks}X_s^{(\lambda)}\Phi,
\]
one has also evidently
\[
 (W_i,W_k)=\sum\limits_s^{1,\ldots,r}c_{iks}W_s\Phi.
\]
It therefore admits $rn-r$ independent solutions
\[
 u_1,u_2,\ldots,u_{rn-r}.
\]
Let us make a change of the independent variables, taking as new variables the
\[
 u_1,u_2,\ldots,u_{rn-r}
\]
together with $r$ new functions independent among themselves and of the $u$, which we shall denote by
\[
 a_1,a_2,\ldots,a_r.
\]
If we transform the $W_i$ into the new variables and denote the transformed operations by $\bar W_i\Phi$, since the $u$ are solutions of (9), we shall have
\[
 \bar W_i\Phi=
 W_i(a_1)\frac{\partial\Phi}{\partial a_1}+
 W_i(a_2)\frac{\partial\Phi}{\partial a_2}+\cdots+
 W_i(a_r)\frac{\partial\Phi}{\partial a_r}.
\]
The new $\bar W_i\Phi$ are still linearly independent and satisfy the relations
\[
 (\bar W_i,\bar W_k)=\sum\limits_s^{1,\ldots,r}c_{iks}\bar W_s\Phi.
\]
But since the derivatives of $\Phi$ with respect to the $u$ do not occur in the $\bar W_i$, it is clear that the property just observed still holds if, for $u_1,u_2,\ldots,u_{rn-r}$, arbitrary constants are put. Calling then
\[
 A_i\Phi=\sum\limits_k^{1,\ldots,r}\alpha_{ik}(a_1,a_2,\ldots,a_r)\frac{\partial\Phi}{\partial a_k}
\]
the expressions thus obtained from the $\bar W_i$, we have evidently constructed the required expressions $A_i\Phi$.

We have thus proved in both parts the second fundamental theorem, which is stated as follows:

Every group $G_r$ with transformations inverse in pairs is generated by $r$ independent infinitesimal transformations
\[
 X_1f,
 X_2f,
 \ldots,
 X_rf,
\]
which satisfy relations of the form
\begin{equation*}\tag{2}
 (X_i,X_j)=\sum\limits_s^{1,\ldots,r} c_{ijs}X_s f,
 \qquad i,j=1,2,\ldots,r,
\end{equation*}
where the $c_{ijs}$ are numerical constants. Conversely, if the $r$ independent infinitesimal transformations $X_if$ satisfy relations of the form (2), they generate such a group $G_r$, which is the totality of all the one-parameter groups $G_1$ generated by the infinitesimal transformations
\[
 \sum\limits_k^{1,\ldots,r} a_kX_kf.
\]

\section*{\S{} 24. The relations among the constants $c_{iks}$}

In the relations (2), which connect the $r$ infinitesimal transformations
\[
 X_1f,
 X_2f,
 \ldots,
 X_rf
\]
generating a group $G_r$, the quantities $c_{ijs}$ are $r^3$ numerical constants. It is now quite important to observe that these $r^3$ constants are not independent, but are connected by certain relations, which we shall now establish.

First of all, since identically
\[
 (X_i,X_k)+(X_k,X_i)=0,
\]
one obtains
\[
 \sum\limits_s^{1,\ldots,r}(c_{iks}+c_{kis})X_s f=0,
\]
and since the $X_sf$ are linearly independent, it follows that
\[
 c_{iks}+c_{kis}=0.
\]
These are the first required relations.

Secondly, if $i,k,j$ are any three indices, Jacobi's identity holds:
\begin{equation*}\tag{11}
 ((X_i,X_k),X_j)+((X_k,X_j),X_i)+((X_j,X_i),X_k)=0.
\end{equation*}
It is proved at once as follows. One has
\[
 ((X_i,X_k),X_j)=X_iX_kX_jf-X_kX_iX_jf-X_jX_iX_kf+X_jX_kX_if,
\]
and since the third term is obtained from the first by circular interchange of the indices, and similarly the fourth from the second, it is clear that by applying twice the cyclic permutation $(ikj)$ and adding the results, one obtains the Jacobian identity (11).

Now in this identity put, for the alternating expressions $(X_i,X_k)$, their expressions as linear functions of the $X$ given by (2). We first find
\[
 ((X_i,X_k),X_j)=
 \left(\sum\limits_s^{1,\ldots,r}c_{iks}X_sf, X_jf\right)
 =\sum\limits_s^{1,\ldots,r}c_{iks}(X_s,X_j)
 =\sum\limits_s^{1,\ldots,r}\sum\limits_a^{1,\ldots,r}c_{iks}c_{sja}X_af,
\]
and therefore (11) becomes
\[
 \sum\limits_a^{1,\ldots,r}\left\{\sum\limits_s^{1,\ldots,r}(c_{iks}c_{sja}+c_{kjs}c_{sia}+c_{jis}c_{ska})\right\}X_af=0.
\]
Since the $X_af$ are independent, it follows that the $c$ must satisfy the equations
\begin{equation*}\tag{I}
\begin{cases}
 c_{iks}+c_{kis}=0,\\[2mm]
 \displaystyle\sum\limits_s^{1,\ldots,r}
  \left(c_{iks}c_{sja}+c_{kjs}c_{sia}+c_{jis}c_{ska}\right)=0,
\end{cases}
\end{equation*}
for all values from $1$ to $r$ of the indices $i,k,j,a$.

We have therefore the important result: the constants $c_{iks}$ of the relations
\[
 (X_i,X_k)=\sum\limits_s^{1,\ldots,r}c_{iks}X_sf
\]
which connect the $r$ infinitesimal transformations $X_1f,X_2f,\ldots,X_rf$ generating a group $G_r$ satisfy equations (I), for all values of the indices from $1$ to $r$. One observes that these relations are entirely independent of the number of variables on which the group operates.

The $r^3$ constants $c_{iks}$ have the greatest importance for the nature of the group, and will become, for a reason that will appear from later studies, the constants of composition or structure of the group.

In order that the $r^3$ constants $c_{iks}$ should belong as constants of composition to a group $G_r$, it is therefore necessary that they satisfy (I). These conditions, as we are now going to prove, are also sufficient: this property is the third fundamental theorem of the theory.

\section*{\S{} 25. Properties of the symbol $\lvert UV\rvert$}

To prove that the relations (I) among the constants $c_{iks}$ are sufficient for there to exist a group $G_r$ which admits the $c_{iks}$ as constants of composition, it will suffice to construct, in a suitable number $m$ of variables, $r$ linearly independent infinitesimal transformations
\[
 X_1f,
 X_2f,
 \ldots,
 X_rf
\]
such that
\[
 (X_i,X_k)=\sum\limits_s^{1,\ldots,r}c_{iks}X_sf.
\]
First consider $r$ independent variables
\[
 y_1,y_2,\ldots,y_r
\]
and, if $U,V$ are any two functions of the $y$, form the expression bilinear in the derivatives of $U,V$ given by
\[
 \sum\limits_{i,k,s}^{1,\ldots,r}c_{iks}y_s
       \frac{\partial U}{\partial y_i}\frac{\partial V}{\partial y_k},
\]
which we denote by the symbol $\lvert UV\rvert$:
\begin{equation*}\tag{12}
 \lvert UV\rvert=
 \sum\limits_{i,k,s}^{1,\ldots,r}c_{iks}y_s
       \frac{\partial U}{\partial y_i}\frac{\partial V}{\partial y_k}.
\end{equation*}
We note that, since
\[
 \lvert U,y_k\rvert=\sum\limits_{i,s}^{1,\ldots,r}c_{iks}y_s\frac{\partial U}{\partial y_i},
\]
one can also write (12) thus:
\begin{equation*}\tag{12*}
 \lvert UV\rvert=\sum\limits_k^{1,\ldots,r}\lvert U,y_k\rvert
                 \frac{\partial V}{\partial y_k}.
\end{equation*}
Now, by the first relation (I) satisfied by the $c_{iks}$, one immediately sees that the symbol $\lvert UV\rvert$ has the first property of alternating expressions, expressed by the formula
\begin{equation*}\tag{13}
 \lvert UV\rvert+\lvert VU\rvert=0.
\end{equation*}
Secondly we say that, if $U,V,W$ are any three functions of the $y$, there also holds the identity analogous to Jacobi's identity
\begin{equation*}\tag{14}
 \lvert\lvert U,V\rvert,W\rvert+
 \lvert\lvert V,W\rvert,U\rvert+
 \lvert\lvert W,U\rvert,V\rvert=0.
\end{equation*}
To this end let us begin by proving that if (14) holds for two given $U,V$ with $W=y_j$, it will also hold with the same $U,V$ and arbitrary $W$. Indeed, by (12*) we have
\[
 \lvert U,W\rvert=\sum\limits_i^{1,\ldots,r}\lvert U,y_i\rvert\frac{\partial W}{\partial y_i}=AW,
 \qquad
 \lvert V,W\rvert=\sum\limits_i^{1,\ldots,r}\lvert V,y_i\rvert\frac{\partial W}{\partial y_i}=BW,
\]
where $AW,BW$ denote the corresponding first-order operations performed on $W$; hence
\[
\begin{split}
&\lvert U,\lvert V,W\rvert\rvert-
 \lvert V,\lvert U,W\rvert\rvert
 =A(BW)-B(AW)                                                    \\
&\quad=\sum\limits_i^{1,\ldots,r}\{A\lvert V,y_i\rvert-B\lvert U,y_i\rvert\}\frac{\partial W}{\partial y_i}
 =\sum\limits_i^{1,\ldots,r}\{\lvert U,\lvert V,y_i\rvert\rvert-
                 \lvert V,\lvert U,y_i\rvert\rvert\}
       \frac{\partial W}{\partial y_i}.
\end{split}
\]
But, (14) holding by hypothesis for $W=y_i$, we have
\[
 \lvert U,\lvert V,y_i\rvert\rvert-
 \lvert V,\lvert U,y_i\rvert\rvert=
 \lvert\lvert U,V\rvert,y_i\rvert,
\]
and therefore the preceding expression becomes
\[
 \lvert U,\lvert V,W\rvert\rvert-
 \lvert V,\lvert U,W\rvert\rvert
 =\sum\limits_i^{1,\ldots,r}\lvert\lvert U,V\rvert,y_i\rvert\frac{\partial W}{\partial y_i}
 =\lvert\lvert U,V\rvert,W\rvert,
\]
which is precisely (14).

It remains only to observe that, putting $U=y_i$, $V=y_k$, $W=y_j$, (14) holds because
\[
 \lvert y_i,y_k\rvert=\sum\limits_s^{1,\ldots,r}c_{iks}y_s
\]
and hence
\[
 \lvert\lvert y_i,y_k\rvert,y_j\rvert
 =\sum\limits_s^{1,\ldots,r} c_{iks}\lvert y_s,y_j\rvert
 =\sum\limits_s^{1,\ldots,r} c_{iks}\sum\limits_a^{1,\ldots,r} c_{sja}y_a;
\]
the identity to be verified follows therefore from the quadratic relations (I) satisfied by the $c_{iks}$.

Thus, since (14) holds for
\[
 U=y_i,
 \quad V=y_k,
 \quad W=y_j,
\]
it also holds for
\[
 U=y_i,
 \quad V=y_k,
 \quad W\hbox{ arbitrary},
\]
then for
\[
 U=y_i,
 \quad V\hbox{ arbitrary},
 \quad W\hbox{ arbitrary},
\]
and finally for arbitrary $U,V,W$. Q. E. D.

If we put
\[
 A_if=\lvert y_i,f\rvert\qquad(i=1,2,\ldots,r),
\]
by the identity (14) we shall have
\[
 (A_i,A_k)=(\lvert y_i,f\rvert,\lvert y_k,f\rvert)=
 \sum\limits_s^{1,\ldots,r} c_{iks}\lvert y_s,f\rvert,
\]
that is,
\[
 (A_i,A_k)=\sum\limits_s^{1,\ldots,r} c_{iks}A_s f.
\]
The $r$ infinitesimal transformations
\[
 A_1f,A_2f,\ldots,A_rf
\]
in the $r$ variables $y_1,y_2,\ldots,y_r$ thus satisfy the required relations of composition. If they are also linearly independent, the goal has already been reached.

But this does not always happen; the simple example in which all the $c_{iks}$ are zero proves it: then all the $A_if$ vanish identically.

\section*{\S{} 26. Change of variables}

To obtain in all cases the required $r$ linearly independent transformations, it is useful first of all to carry out a certain change of independent variables, from $y_1,y_2,\ldots,y_r$ to $Y_1,Y_2,\ldots,Y_r$, such that all the symbols $\lvert Y_i,Y_k\rvert$ have the value zero or one. This is obtained as follows.

First, if all the $c_{iks}$ are zero, it is enough to take $Y_i=y_i$; otherwise put, for example, that not all the $c_{1ks}$ are zero and take a solution $X_1$ of the partial differential equation
\[
 \lvert y_1,f\rvert=\sum\limits_{k,s}^{1,\ldots,r}c_{1ks}y_s\frac{\partial f}{\partial y_k}=1,
\]
and, putting $P_1=y_1$, the pair of functions $P_1,X_1$ satisfies the condition
\[
 \lvert P_1,X_1\rvert=1.
\]
Consider now the two simultaneous equations
\begin{equation*}\tag{15}
 A_1f=\lvert P_1,f\rvert=0,
 \qquad
 A_2f=\lvert X_1,f\rvert=0,
\end{equation*}
which are certainly independent because $f=P_1$ satisfies the first and not the second. They also form a complete system because
\[
 (A_1,A_2)=\lvert P_1,\lvert X_1,f\rvert\rvert-
 \lvert X_1,\lvert P_1,f\rvert\rvert=
 \lvert\lvert P_1,X_1\rvert,f\rvert=
 \lvert1,f\rvert=0;
\]
they therefore admit $r-2$ independent solutions
\[
 y'_1,y'_2,\ldots,y'_{r-2}.
\]
These are also independent of $P_1,X_1$, because, if they were connected by a relation
\[
 \varphi(P_1,X_1,y'_1,y'_2,\ldots,y'_{r-2})=0,
\]
the operations $A_1,A_2$ performed on $\varphi$ would give
\[
 \frac{\partial\varphi}{\partial X_1}=0,
 \qquad
 \frac{\partial\varphi}{\partial P_1}=0.
\]
Moreover, applying identity (14) with
\[
 U=y'_i,
 \quad V=y'_k,
 \quad W=\begin{cases}P_1,\\ X_1,
 \end{cases}
\]
we see that $\lvert y'_i,y'_k\rvert$ is also a solution of system (15), and therefore is a function of $y'_1,y'_2,\ldots,y'_{r-2}$ alone. Thus the $r$ independent functions
\[
 P_1,
 \quad X_1,
 \quad y'_1,y'_2,\ldots,y'_{r-2}
\]
satisfy the conditions
\[
 \lvert P_1,X_1\rvert=1,
 \quad \lvert P_1,y'_i\rvert=0,
 \quad \lvert X_1,y'_i\rvert=0,
 \quad
 \lvert y'_i,y'_k\rvert=\varphi_{ik}(y'_1,y'_2,\ldots,y'_{r-2}).
\]
If all the $\varphi_{ik}$ were zero, we would already have obtained the required change of variables. Otherwise, proceeding with
\[
 y'_1,y'_2,\ldots,y'_{r-2}
\]
as before with $y_1,y_2,\ldots,y_r$, we shall find $r-2$ independent functions of the $y'$
\[
 P_2,
 X_2,
 y''_1,y''_2,\ldots,y''_{r-4}
\]
such that
\[
 \lvert P_2,X_2\rvert=1,
 \quad \lvert P_2,y''_i\rvert=0,
 \quad \lvert X_2,y''_i\rvert=0,
 \quad
 \lvert y''_i,y''_k\rvert=\varphi''_{ik}(y''_1,y''_2,\ldots,y''_{r-4}).
\]
Continuing in this way, we shall find $r=2m+q$ independent functions of $y_1,y_2,\ldots,y_r$
\[
 P_1,P_2,\ldots,P_m,
 \quad X_1,X_2,\ldots,X_m,
 \quad Y_1,Y_2,\ldots,Y_q
\]
such that
\begin{equation*}\tag{16}
 \lvert P_i,P_k\rvert=0,
 \quad \lvert P_i,X_k\rvert=\varepsilon_{ik},
 \quad \lvert X_i,X_k\rvert=0,
 \quad \lvert Y_i,Y_k\rvert=0.
\end{equation*}
Taking the $P$ and the $X$ as new variables, the proposed goal is attained.

\section*{\S{} 27. The third fundamental theorem}

Expressing the symbol $\lvert UV\rvert$ in the new variables $P_i,X_i$ we shall have
\[
\begin{split}
\lvert UV\rvert={}&
 \sum\limits_{i,k}^{1,\ldots,m+q}\lvert P_i,P_k\rvert
        \frac{\partial U}{\partial P_i}\frac{\partial V}{\partial P_k}
 +\sum\limits_{i,k}^{1,\ldots,m}\lvert X_i,X_k\rvert
        \frac{\partial U}{\partial X_i}\frac{\partial V}{\partial X_k}        \\
&+\sum\limits_i^{1,\ldots,m+q}\sum\limits_k^{1,\ldots,m}
 \lvert P_i,X_k\rvert
 \left(
  \frac{\partial U}{\partial P_i}\frac{\partial V}{\partial X_k}
 -\frac{\partial U}{\partial X_k}\frac{\partial V}{\partial P_i}
 \right).
\end{split}
\]
And therefore, by (16),
\[
 \lvert UV\rvert=
 \sum\limits_i^{1,\ldots,m}
 \left(
  \frac{\partial U}{\partial P_i}\frac{\partial V}{\partial X_i}
 -\frac{\partial U}{\partial X_i}\frac{\partial V}{\partial P_i}
 \right).
\]
Denoting by
\[
 \Omega_k=\Omega_k(P,X)
\]
the transformation formulae of the variables, the $r$ operations
\[
 A_k f=\lvert y_k,f\rvert\qquad (k=1,2,\ldots,r)
\]
constructed in \S{}25, expressed in the new variables, become
\begin{equation*}\tag{17}
 A_kf=\sum\limits_i^{1,\ldots,m}
 \left(
  \frac{\partial\Omega_k}{\partial P_i}\frac{\partial f}{\partial X_i}
 -\frac{\partial\Omega_k}{\partial X_i}\frac{\partial f}{\partial P_i}
 \right).
\end{equation*}

Naturally, also in the new variables the $A_i f$ satisfy the relations
\[
 (A_i,A_k)=\sum\limits_s^{1,\ldots,r} c_{iks}A_s f;
\]
however, the number of linearly independent ones among the $A_i f$ has still remained the same, and up to now we have gained no advantage over the primitive treatment. On the other hand, the form of (17) suggests a simple device for attaining the goal. To the $r=2m+q$ variables
\[
 P_1,P_2,\ldots,P_{m+q};\quad X_1,X_2,\ldots,X_m,
\]
we add $q$ new independent variables
\[
 X_{m+1},X_{m+2},\ldots,X_{m+q}
\]
and form
\begin{equation*}\tag{18}
 B_k f=\sum\limits_i^{1,\ldots,m+q}
 \left(
  \frac{\partial\Omega_k}{\partial P_i}\frac{\partial f}{\partial X_i}
 -\frac{\partial\Omega_k}{\partial X_i}\frac{\partial f}{\partial P_i}
 \right),
\end{equation*}
that is, since the $\Omega$ do not contain $X_{m+1},\ldots,X_{m+q}$,
\begin{equation*}\tag{18*}
 B_k f=A_k f+\sum\limits_{\alpha}^{1,\ldots,q}
 \frac{\partial\Omega_k}{\partial P_{m+\alpha}}
 \frac{\partial f}{\partial X_{m+\alpha}}.
\end{equation*}
I say that the $r$ infinitesimal transformations
\[
 B_1f,\ B_2f,\ldots,\ B_rf
\]
in the $2m+2q$ variables $P,X$ satisfy the composition relations
\[
 (B_i,B_k)=\sum\limits_s^{1,\ldots,r} c_{iks}B_sf
\]
and are linearly independent.

The first assertion follows by observing that from (18*), since $X_{m+1},\ldots,X_{m+q}$ do not occur in the coefficients of the $A_k$, one has
\[
 (B_i,B_k)=(A_i,A_k)+\sum\limits_{\alpha}^{1,\ldots,q}
 \left\{A_i\left(\frac{\partial\Omega_k}{\partial P_{m+\alpha}}\right)
       -A_k\left(\frac{\partial\Omega_i}{\partial P_{m+\alpha}}\right)
 \right\}
 \frac{\partial f}{\partial X_{m+\alpha}}.
\]
But one finds immediately from (17)
\[
\begin{split}
&A_i\left(\frac{\partial\Omega_k}{\partial P_{m+\alpha}}\right)
 -A_k\left(\frac{\partial\Omega_i}{\partial P_{m+\alpha}}\right)       \\
&\quad=\frac{\partial}{\partial P_{m+\alpha}}
 \sum\limits_{\nu}^{1,\ldots,m}
 \left(
  \frac{\partial\Omega_i}{\partial P_\nu}
  \frac{\partial\Omega_k}{\partial X_\nu}
 -\frac{\partial\Omega_i}{\partial X_\nu}
  \frac{\partial\Omega_k}{\partial P_\nu}
 \right)                                                   \\
&\quad=\frac{\partial}{\partial P_{m+\alpha}}\lvert y_i,y_k\rvert
 =\frac{\partial}{\partial P_{m+\alpha}}\sum\limits_s^{1,\ldots,r}c_{iks}\Omega_s
 =\sum\limits_s^{1,\ldots,r} c_{iks}\frac{\partial\Omega_s}{\partial P_{m+\alpha}}.
\end{split}
\]
Consequently
\[
\begin{split}
 (B_i,B_k)&=\sum\limits_s^{1,\ldots,r}c_{iks}A_s f
 +\sum\limits_s^{1,\ldots,r}c_{iks}\sum\limits_{\alpha}^{1,\ldots,q}
 \frac{\partial\Omega_s}{\partial P_{m+\alpha}}
 \frac{\partial f}{\partial X_{m+\alpha}}                 \\
 &=\sum\limits_s^{1,\ldots,r}c_{iks}B_s f,
\end{split}
\]
as was required.

Now suppose that among the $B_if$ there holds a relation
\[
 \sum\limits_i^{1,\ldots,r} e_iB_i f=0
\]
with constant $e$. Taking in it successively
\[
 f=P_1,P_2,\ldots,P_{m+q};\quad X_1,X_2,\ldots,X_{m+q},
\]
one would obtain
\[
 \sum\limits_i^{1,\ldots,r} e_i\frac{\partial\Omega_i}{\partial X_\lambda}=0,
 \qquad
 \sum\limits_i^{1,\ldots,r} e_i\frac{\partial\Omega_i}{\partial P_\lambda}=0,
 \qquad \lambda=1,2,\ldots,m+q,
\]
and hence
\[
 \sum\limits_i^{1,\ldots,r} e_i\Omega_i=\sum\limits_i^{1,\ldots,r} e_iy_i=\hbox{constant},
\]
which, since the $y$ are independent, gives
\[
 e_1=e_2=\cdots=e_r=0.
\]
Thus the required $r$ transformations have in fact been constructed, and we may finally formulate the fundamental theorem of the theory of groups, stated as follows.

If a group $G_r$, in any number $n$ of variables $x_1,x_2,\ldots,x_n$, possesses the $r$ infinitesimal transformations
\[
 X_i f=\sum\limits_k^{1,\ldots,n}\xi_{ik}(x_1,x_2,\ldots,x_n)\frac{\partial f}{\partial x_k},
\]
the constants $c_{iks}$ of the composition relations
\[
 (X_i,X_k)=\sum\limits_s^{1,\ldots,r} c_{iks}X_s f
\]
satisfy the relations
\[
\left\{
\begin{array}{l}
 c_{iks}+c_{kis}=0,\\[0.4em]
 \displaystyle\sum\limits_s^{1,\ldots,r}
 \left(c_{iks}c_{sja}+c_{kjs}c_{sia}+c_{jis}c_{ska}\right)=0.
\end{array}
\right.
\tag{I}
\]
Conversely, if the $r^3$ constants $c_{iks}$ satisfy (I), one can always find, in a suitable number of variables, $r$ linearly independent infinitesimal transformations for which the composition constants have the values $c_{iks}$.

\section*{\S{} 28. The Maurer equations for the $\psi_{ik}$}

The three theorems that we have proved in the present chapter are of fundamental importance for the theory of groups, and it is useful to know different demonstrations of them. Already in \S{}22 we gave a new proof of the first theorem, and now we propose to do the same for the other two. The new proofs are again due to Schur, and we shall present them with the modifications and simplifications introduced by Engel and by Lie himself.

We shall therefore prove that if the $r^3$ constants $c_{ijs}$ satisfy the relations (I), then one can find, in $r$ variables
\[
 a_1,a_2,\ldots,a_r,
\]
$r$ linearly independent transformations
\begin{equation*}\tag{19}
 A_i f=\sum\limits_k^{1,\ldots,r}\alpha'_{ik}(a_1,a_2,\ldots,a_r)
        \frac{\partial f}{\partial a_k}
\end{equation*}
which satisfy the composition relations
\begin{equation*}\tag{20}
 (A_i,A_j)=\sum\limits_s^{1,\ldots,r}c_{ijs}A_sf.
\end{equation*}
In this way the third fundamental theorem will be proved anew, and at the same time also the second, as follows immediately from \S{}21.

It is thus a matter of finding $r^2$ functions $\alpha'_{ik}$ of $a_1,a_2,\ldots,a_r$, whose determinant is $\ne0$--which assures the linear independence of (19)--and which satisfy the relations (20), namely
\[
 A_i(\alpha'_{j\mu})-A_j(\alpha'_{i\mu})=\sum\limits_s^{1,\ldots,r} c_{ijs}\alpha'_{s\mu},
\]
or equivalently the following
\begin{equation*}\tag{21}
 \sum\limits_k^{1,\ldots,r}
 \left(\alpha'_{ik}\frac{\partial\alpha'_{j\mu}}{\partial a_k}
      -\alpha'_{jk}\frac{\partial\alpha'_{i\mu}}{\partial a_k}\right)
 =\sum\limits_s^{1,\ldots,r} c_{ijs}\alpha'_{s\mu}.
\end{equation*}
In agreement with the notation already used in \S{}19, let $\psi_{hk}$ denote the quotient of the algebraic complement of $\alpha'_{kh}$ in the determinant of the $\alpha'$ by that determinant itself. The equations (21) can be changed into equivalent equations for the $\psi$ in the following way.

Multiply (21) by $\psi_{\nu\mu}$, where $i,j,\nu$ are arbitrary indices from $1$ to $r$, and sum with respect to $\mu$; we obtain
\begin{equation*}\tag{22}
 \sum\limits_{k,\mu}^{1,\ldots,r}
 \left(\alpha'_{ik}\psi_{\nu\mu}\frac{\partial\alpha'_{j\mu}}{\partial a_k}
      -\alpha'_{jk}\psi_{\nu\mu}\frac{\partial\alpha'_{i\mu}}{\partial a_k}\right)
 =c_{ij\nu}.
\end{equation*}
But one has
\[
 \sum\limits_\mu^{1,\ldots,r}\psi_{\nu\mu}\alpha'_{j\mu}=\varepsilon_{\nu j},
 \qquad
 \sum\limits_\mu^{1,\ldots,r}\psi_{\nu\mu}\alpha'_{i\mu}=\varepsilon_{\nu i},
\]
and hence, differentiating,
\[
 \sum\limits_\mu^{1,\ldots,r}\psi_{\nu\mu}\frac{\partial\alpha'_{j\mu}}{\partial a_k}
 =-\sum\limits_\mu^{1,\ldots,r}\alpha'_{j\mu}\frac{\partial\psi_{\nu\mu}}{\partial a_k},
\]
\[
 \sum\limits_\mu^{1,\ldots,r}\psi_{\nu\mu}\frac{\partial\alpha'_{i\mu}}{\partial a_k}
 =-\sum\limits_\mu^{1,\ldots,r}\alpha'_{i\mu}\frac{\partial\psi_{\nu\mu}}{\partial a_k}.
\]
Thus (22) becomes
\[
 \sum\limits_{k,\mu}^{1,\ldots,r}\alpha'_{jk}\alpha'_{i\mu}\frac{\partial\psi_{\nu\mu}}{\partial a_k}
 -\sum\limits_{k,\mu}^{1,\ldots,r}\alpha'_{ik}\alpha'_{j\mu}\frac{\partial\psi_{\nu\mu}}{\partial a_k}=c_{ij\nu};
\]
and, permuting in the second sum the summation indices $k,\mu$, we have
\[
 \sum\limits_{k,\mu}^{1,\ldots,r}\alpha'_{jk}\alpha'_{i\mu}
 \left(\frac{\partial\psi_{\nu\mu}}{\partial a_k}
      -\frac{\partial\psi_{\nu k}}{\partial a_\mu}\right)=c_{ij\nu}.
\]
Finally, multiplying this last equation by $\psi_{i\lambda}\psi_{j\nu}$ and summing with respect to $i,j$, one obtains the required equations
\begin{equation*}\tag{II}
 \frac{\partial\psi_{k\lambda}}{\partial a_\nu}
 -\frac{\partial\psi_{k\nu}}{\partial a_\lambda}
 =\sum\limits_{i,j}^{1,\ldots,r}c_{ijk}\psi_{i\lambda}\psi_{j\nu}.
\end{equation*}
As the (II) are deduced from (22), so, by reversing the calculation just carried out, one sees that (22) follows from (II), to which therefore they are perfectly equivalent.

The equations (II) were established by Maurer. The proposed question is therefore reduced to finding $r^2$ functions $\psi_{ik}$ of $a_1,a_2,\ldots,a_r$, with non-zero determinant, and satisfying Maurer's equations (II).

\section*{\S{} 29. Schur's proof of the second and third fundamental theorems}

By a suitable device one can prove, as Schur did, a system of $r^2$ functions $\psi$ that satisfy (II), integrating only a system of ordinary linear differential equations with constant coefficients; this requires only, as is well known, algebraic operations, namely the solution of the so-called characteristic equation. Thus not only are the required infinitesimal transformations
\[
 A_1f,A_2f,\ldots,A_rf,
\]
proved to exist, but it is shown that only algebraic operations are needed for this.

For this purpose we perform in (II) a change of variables and of unknown functions, putting
\[
 a_1=\lambda_1t,
 \quad a_2=\lambda_2t,
 \quad \ldots,
 \quad a_r=\lambda_rt,
\]
\[
 \psi_{ik}(a_1,a_2,\ldots,a_r)=\theta_{ik}(\lambda_1,\lambda_2,\ldots,\lambda_r;t);
\]
with this, the equations (II)
\[
 \frac{\partial\psi_{s\alpha}}{\partial a_\beta}
 -\frac{\partial\psi_{s\beta}}{\partial a_\alpha}
 =\sum\limits_{i,j}^{1,\ldots,r}c_{ijs}\psi_{i\alpha}\psi_{j\beta}
\]
become
\begin{equation*}\tag{II*}
 \frac{\partial\theta_{s\alpha}}{\partial\lambda_\beta}
 -\frac{\partial\theta_{s\beta}}{\partial\lambda_\alpha}
 =\sum\limits_{i,j}^{1,\ldots,r}c_{ijs}\theta_{i\alpha}\theta_{j\beta}.
\end{equation*}
Now integrate the system of ordinary linear differential equations with respect to $t$
\begin{equation*}\tag{23}
 \frac{\partial\theta_{s\alpha}}{\partial t}
 =\varepsilon_{s\alpha}+\sum\limits_{i,j}^{1,\ldots,r}c_{ijs}\lambda_j\theta_{i\alpha},
 \qquad s,\alpha=1,2,\ldots,r,
\end{equation*}
with the initial conditions
\begin{equation*}\tag{24}
 \theta_{s\alpha}=0\qquad\hbox{for }t=0,
\end{equation*}
where, in the integration, the $\lambda$ are to be regarded as constants (parameters). From the form of these linear differential equations with constant coefficients it follows immediately that the functions $\theta_{s\alpha}$ will contain the variables $\lambda_1,\lambda_2,\ldots,\lambda_r,t$ precisely in the products
\[
 \lambda_1t,\quad \lambda_2t,\quad \ldots,
 \quad \lambda_rt;
\]
and we shall now prove that the $\theta_{s\alpha}$ so determined satisfy (II*), and therefore that the $\psi_{s\alpha}$ satisfy (II), as was required.

Indeed put
\begin{equation*}\tag{25}
 V_{s\alpha\beta}=\frac{\partial\theta_{s\alpha}}{\partial\lambda_\beta}
 -\frac{\partial\theta_{s\beta}}{\partial\lambda_\alpha}
 -\sum\limits_{i,j}^{1,\ldots,r}c_{ijs}\theta_{i\alpha}\theta_{j\beta}.
\end{equation*}
Differentiating with respect to $t$ and using (23), one obtains
\[
\begin{split}
\frac{\partial V_{s\alpha\beta}}{\partial t}={}&
\frac{\partial}{\partial\lambda_\beta}
 \left[\varepsilon_{s\alpha}+\sum\limits_{i,j}^{1,\ldots,r}c_{ijs}\lambda_j\theta_{i\alpha}\right]
-\frac{\partial}{\partial\lambda_\alpha}
 \left[\varepsilon_{s\beta}+\sum\limits_{i,j}^{1,\ldots,r}c_{ijs}\lambda_j\theta_{i\beta}\right]                 \\
&-\sum\limits_{i,j}^{1,\ldots,r}c_{ijs}\theta_{j\beta}
 \left[\varepsilon_{i\alpha}+\sum\limits_{k,v}^{1,\ldots,r}c_{kvi}\lambda_v\theta_{k\alpha}\right]
-\sum\limits_{i,j}^{1,\ldots,r}c_{ijs}\theta_{i\alpha}
 \left[\varepsilon_{j\beta}+\sum\limits_{k,v}^{1,\ldots,r}c_{kvj}\lambda_v\theta_{k\beta}\right].
\end{split}
\]
That is,
\[
\begin{split}
\frac{\partial V_{s\alpha\beta}}{\partial t}={}&
 \sum\limits_i^{1,\ldots,r} c_{i\beta s}\theta_{i\alpha}
-\sum\limits_i^{1,\ldots,r} c_{i\alpha s}\theta_{i\beta}
+\sum\limits_{i,j}^{1,\ldots,r}c_{ijs}\lambda_j
 \left(\frac{\partial\theta_{i\alpha}}{\partial\lambda_\beta}
      -\frac{\partial\theta_{i\beta}}{\partial\lambda_\alpha}\right)                         \\
&-\sum\limits_{i,j}^{1,\ldots,r}c_{ijs}\theta_{j\beta}\varepsilon_{i\alpha}
-\sum\limits_{i,j}^{1,\ldots,r}c_{ijs}\theta_{i\alpha}\varepsilon_{j\beta} \\
&-\sum\limits_{i,j,k,v}^{1,\ldots,r}c_{ijs}c_{kvi}\lambda_v\theta_{j\beta}\theta_{k\alpha}
-\sum\limits_{i,j,k,v}^{1,\ldots,r}c_{ijs}c_{kvj}\lambda_v\theta_{i\alpha}\theta_{k\beta}.
\end{split}
\]
The simple sums cancel by the first of the relations (I) for the $c$, and, putting for
\[
 \frac{\partial\theta_{i\alpha}}{\partial\lambda_\beta}
 -\frac{\partial\theta_{i\beta}}{\partial\lambda_\alpha}
\]
the expression given by (25), there remains
\[
\begin{split}
\frac{\partial V_{s\alpha\beta}}{\partial t}={}&
 \sum\limits_{i,j}^{1,\ldots,r}c_{ijs}\lambda_jV_{i\alpha\beta}
-\sum\limits_{i,j,k,v}^{1,\ldots,r}c_{ijs}c_{kvi}\lambda_v\theta_{j\beta}\theta_{k\alpha}  \\
&-\sum\limits_{i,j,k,v}^{1,\ldots,r}c_{ijs}c_{kvj}\lambda_v\theta_{i\alpha}\theta_{k\beta}.
\end{split}
\]
By permuting in the second of the quadruple sums the indices $j,v$, and in the last carrying out the cyclic substitution $(ikv)$ on the indices, one sees that the two quadruple sums cancel. Thus the $V_{s\alpha\beta}$ defined by (25) satisfy the homogeneous linear equations
\[
 \frac{\partial V_{s\alpha\beta}}{\partial t}
 =\sum\limits_{i,j}^{1,\ldots,r}c_{ijs}\lambda_jV_{i\alpha\beta}.
\]
On the other hand, for $t=0$ they vanish, by reason of the initial conditions (24), and consequently they are always zero; this proves that (II*) holds. Thus we have proved:

Given a system of $r^3$ constants $c_{iks}$ which verify (I), one can find, by algebraic operations alone, $r$ independent infinitesimal transformations
\[
 A_i f=\sum\limits_k^{1,\ldots,r}\alpha_{ik}(a_1,a_2,\ldots,a_r)\frac{\partial f}{\partial a_k}
\]
in $r$ variables $a_1,a_2,\ldots,a_r$ which satisfy the composition relations
\[
 (A_i,A_k)=\sum\limits_s^{1,\ldots,r} c_{iks}A_s f
\]
and therefore generate a group $G_r$.

\section*{\S{} 30. Composition of two $\infty^r$ series of transformations}

With the three theorems with which we have occupied ourselves in the present chapter, we have laid the foundations for the general theory of continuous finite groups $G_r$ with transformations inverse in pairs, whose finite equations are contained in a single $\infty^r$ series
\[
 x'_i=f_i(x_1,x_2,\ldots,x_n;a_1,a_2,\ldots,a_r).
\]
Already in \S{}4 we indicated that one may consider more general groups with $r$ parameters whose transformations, instead of forming a single $\infty^r$ series, constitute several such series or even an infinite but discrete number of them. We now wish to prove what was asserted then, namely that the study of such mixed groups is reduced to groups $G_r$ combined with ordinary discontinuous groups.

For this it will be useful to take up again, extending them, the considerations of \S\S{}12--14; and first of all we shall consider the following question.

Given two series, each $\infty^r$, of transformations
\begin{align*}
 x'_i&=f_i(x_1,x_2,\ldots,x_n;a_1,a_2,\ldots,a_r), \tag{26}\\
 x''_i&=\varphi_i(x'_1,x'_2,\ldots,x'_n;b_1,b_2,\ldots,b_r), \tag{27}
\end{align*}
if we compose any transformation $S_a$ of the first with any $T_b$ of the second, in the composite transformation $S_aT_b$
\[
 (S_aT_b)\quad
 x''_i=\varphi_i(f_1(x,a),\ldots,f_n(x,a);b_1,\ldots,b_r)
\]
the number of essential parameters is at most $2r$ and at least $r$\begingroup\renewcommand\thefootnote{(*)}\footnotemark\endgroup.\begingroup\renewcommand\thefootnote{(*)}\footnotetext{This is evident because, keeping for example the $b$ fixed and letting the $a$ vary, $S_aT_b$ runs through $\infty^r$ transformations.}\endgroup{} We ask: what relations must hold between the two $\infty^r$ series $S_a,T_b$ in order that the composed series $S_aT_b$ still have only $r$ essential parameters?

Denoting by $c_1,c_2,\ldots,c_r$ the $r$ essential parameters of the composed series, we may write
\begin{equation*}\tag{28}
 x''_i=\psi_i(x_1,x_2,\ldots,x_n;c_1,c_2,\ldots,c_r),
\end{equation*}
where the parameters $c$ are functions of the $a$ and the $b$
\begin{equation*}\tag{29}
 c_\eta=\theta_\eta(a_1,\ldots,a_r;b_1,\ldots,b_r).
\end{equation*}
From the same observation made at the beginning of \S{}13 it follows that these relations (28) are solvable both with respect to the $a$ and with respect to the $b$, that is, the functional determinants
\[
 \frac{\partial(\theta_1,\theta_2,\ldots,\theta_r)}{\partial(a_1,a_2,\ldots,a_r)},
 \qquad
 \frac{\partial(\theta_1,\theta_2,\ldots,\theta_r)}{\partial(b_1,b_2,\ldots,b_r)}
\]
will each be different from zero. Now give to the $a,b$ infinitesimal increments $\delta a,\delta b$ such as to satisfy the conditions $\delta c_\eta=0$, that is
\begin{equation*}\tag{30}
 \sum\limits_k^{1,\ldots,r}\frac{\partial\theta_\eta}{\partial a_k}\delta a_k
 +\sum\limits_k^{1,\ldots,r}\frac{\partial\theta_\eta}{\partial b_k}\delta b_k=0,
\end{equation*}
in which, from what was said above, one may choose the $\delta a$ arbitrarily and calculate the $\delta b$ inversely. Precisely as in \S{}13 there will follow
\begin{equation*}\tag{31}
 S_a^{-1}S_{a+\delta a}=T_bT_{b+\delta b}^{-1}.
\end{equation*}
Fixing the $a$ arbitrarily and letting the $\delta a$ vary arbitrarily, the first member of (31) runs through an $\infty^{r-1}$ series of infinitesimal transformations; the same is true of the second member when the $b$ are fixed and the $\delta b$ are made to vary arbitrarily.

The equality (31) therefore expresses (cf. \S{}13) that these two $\infty^{r-1}$ series of infinitesimal transformations always remain the same, whatever values are assigned to the $a$ or to the $b$.

Therefore, since the transformations of the series (26) infinitesimally near to $S_a$ are obtained by combining $S_a$ with an infinitesimal transformation of the aforesaid fixed series, the reasoning of \S{}14 is again applicable, and it follows that the series (26) satisfies a system of differential equations of the type
\[
 \frac{\partial x'_i}{\partial a_k}=\sum\limits_j^{1,\ldots,r}\eta_{jk}(a)\xi_{ji}(x'),
\]
for which the determinant of the $\eta$ is not identically zero. If, then, we give to the $a$ a particular system of values $a^{(0)}$, such only that
\[
 \lvert\eta_{jk}(a^{(0)})\rvert\ne0,
\]
it follows from the theorem at the end of \S{}16 that the transformations
\[
 S_{a^{(0)}}^{-1}S_a
\]
form the group $G_r$, with transformations inverse in pairs, generated by the $r$ infinitesimal transformations
\[
 X_j f=\sum\limits_i^{1,\ldots,n}\xi_{ji}(x_1,x_2,\ldots,x_n)\frac{\partial f}{\partial x_i}.
\]
On the other hand, if we fix any transformation $T_{b^{(0)}}$ of the series (27), we can satisfy the relation
\[
 S_aT_b=S_{a^{(0)}}T_{b^{(0)}},
\]
that is, the relations
\begin{equation*}\tag{32}
 \theta_\lambda(a,b)=\theta_\lambda(a^{(0)},b^{(0)})
 \qquad (\lambda=1,2,\ldots,r),
\end{equation*}
by taking $S_a$ arbitrarily and calculating $T_b$ inversely. If we write (32) in the form
\[
 S_{a^{(0)}}^{-1}S_a=T_{b^{(0)}}T_b^{-1},
\]
we see that $T_{b^{(0)}}T_b^{-1}$ runs through the same transformations as $S_{a^{(0)}}^{-1}S_a$, that is, the transformations of the group $G_r$. Denoting these latter by the symbol $E_\lambda$, we shall have
\[
 S_a=S_{a^{(0)}}E_\lambda,
\]
and since the inverse $T_bT_{b^{(0)}}^{-1}$ of $T_{b^{(0)}}T_b^{-1}$ is also an $E$,
\[
 T_b=E_\mu T_{b^{(0)}}.
\]
Conversely, if $E_\lambda$ runs through the transformations of a group $G_r$ and $S_{a^{(0)}},T_{b^{(0)}}$ are two arbitrary fixed transformations, it is clear that the two $\infty^r$ series
\[
 S_{a^{(0)}}E_\lambda,
 \qquad
 E_\lambda T_{b^{(0)}}
\]
have the required property, because
\[
 (S_{a^{(0)}}E_\lambda)(E_\mu T_{b^{(0)}})
 =S_{a^{(0)}}(E_\lambda E_\mu)T_{b^{(0)}}
 =S_{a^{(0)}}E_\nu T_{b^{(0)}}
\]
contains only the $r$ parameters of the composite transformation $E_\nu$. We have therefore established the theorem:

In order that two $\infty^r$ series of transformations $S_a,T_b$ give rise to a composite series $S_aT_b$ with only $r$ essential parameters, it is necessary and sufficient that there exist a group $G_r$, with transformations inverse in pairs, such that, denoting by $E_\lambda$ a variable transformation in $G_r$ and by $A,B$ two fixed transformations, one of the first and the other of the second series, every transformation of the first series be given by
\[
 S_a=AE_\lambda
\]
and every one of the second by
\[
 T_b=E_\mu B.
\]
If in particular we suppose that the two $\infty^r$ series coincide, one must have
\[
 S_{a^{(0)}}E_\lambda=E_\mu S_{a^{(0)}},
\]
that is
\[
 S_{a^{(0)}}^{-1}E_\mu S_{a^{(0)}}=E_\lambda,
\]
whence we have the theorem:

If an $\infty^r$ series of transformations $S_a$ has the property that the composite transformation of any two of the series has only $r$ essential parameters, the series itself has the form
\[
 S_a\equiv AE_a,
\]
where $E_a$ runs through the transformations of a group $G_r$ and $A$ is any fixed transformation in the series, which must satisfy the condition
\[
 A^{-1}G_rA=G_r,
\]
that is, it will be permutable with the group.

\section*{\S{} 31. Series $\infty^r$ of transformations in which products of $m$ transformations are in the series}

We now pass to the solution of the following question. What conditions must an $\infty^r$ series of transformations satisfy so that, on composing any $m$ transformations of the series, a transformation of the same series always results, while this does not occur when a smaller number is composed?

We first observe that the composite transformation of two transformations of the series must contain only $r$ parameters; for if it contained more, these could no longer disappear by composing with the succeeding ones. Applying the second theorem of the preceding paragraph, we see meanwhile that the transformations of our series will certainly have the form
\[
 S_a=AE_a
\]
of that theorem. Now, composing any two transformations of the series
\[
 AE_\lambda,
 \qquad
 AE_\mu,
\]
since $AE_\mu A^{-1}=E_\nu$ and $AE_\mu=E_\nu A$, one has
\[
 (AE_\lambda)(AE_\mu)=A^2E_\nu.
\]
In general, by composing $s$ transformations of the series, one will have a transformation of the form
\[
 A^sE_a.
\]
Our hypothesis therefore entails that one has
\[
 A^mE_\lambda=AE_\mu,
\]
that is
\[
 A^{m-1}=E_\nu,
\]
which is to say that the $(m-1)$-st power of $A$ belongs to $G_r$, but no other power of lower order. We conclude that every series of the desired kind is obtained with the substitutions $E_\lambda$ of a group $G_r$ with transformations inverse in pairs, a fixed transformation $A$ permutable with the group $G_r$, and such that the power $A^{m-1}$ of $A$, but

no power of lower order, belongs to $G_r$.

It is necessary to recall that, for the exact application of this theorem, it is proper to limit ourselves to considering those transformations of the series which lie in a suitably restricted neighbourhood of $A$, since the validity of the theorem at the end of \S{}16, on which we have relied, is subject precisely to such a limitation.

If the preceding theorem is applied to the case $m=2$, that is, under the hypothesis that the series $S_a$ constitutes a group $\Gamma$, without supposing that in $\Gamma$ there is the identity or pairs of inverse transformations, we see that in that case the transformations of $\Gamma$ belong to a group $G_r$ with transformations inverse in pairs. It follows that the equations of any finite group with $r$ parameters are obtained from those of a group $G_r$ with transformations inverse in pairs by a change of parameters and by analytic continuation. Thus there appears the appropriateness, in the theory of finite continuous groups, of limiting oneself to the consideration of the groups $G_r$ with transformations inverse in pairs.

\section*{\S{} 32. Mixed groups}

Suppose now that we have a number, finite at first, say $m$, of series of transformations, each containing $r$ parameters,
\[
 x'_i=f_i^{(k)}(x_1,x_2,\ldots,x_n; a_1^{(k)},a_2^{(k)},\ldots,a_r^{(k)}),
 \qquad
 \begin{gathered}
 i=1,2,\ldots,n,\\
 k=1,2,\ldots,m,
 \end{gathered}
\]
and that these $m$ series, considered together, form a mixed group. Considering each of these series by itself and comparing them two by two, the theorems of \S{}30 will be applicable. Hence, denoting by $S^1,S^2,\ldots,S^m$ $m$ fixed transformations taken in the respective series, and by
\[
 S_{a^{(1)}}^1,
 \quad S_{a^{(2)}}^2,
 \quad \ldots,
 \quad S_{a^{(m)}}^m
\]
variable transformations in the same series, it follows that the $m$ series
\[
 (S^1)^{-1}S_{a^{(1)}}^1,
 \quad (S^2)^{-1}S_{a^{(2)}}^2,
 \quad \ldots,
 \quad (S^m)^{-1}S_{a^{(m)}}^m
\]
coincide in a single group $G_r$ with transformations inverse in pairs, a group permutable with all the transformations of the series. Denoting by $E_a$ a variable transformation in $G_r$, the $m$ series will be given by the transformations
\begin{equation*}\tag{33}
 S^1E_a,
 \quad S^2E_a,
 \quad \ldots,
 \quad S^mE_a.
\end{equation*}
Now, if all these transformations are multiplied by any one of them, only transformations forming a group will be reproduced; moreover, since the $S$ are permutable with the group $G_r$, the $m$ series will do nothing but permute among themselves. Conversely, the product $S_iS_k$, belonging to the same series, necessarily has the form
\[
 S_iS_k=S_\ell E_{\lambda'},
\]
and one has
\[
\begin{split}
 (S_iE_\lambda)(S_kE_\mu)
  &=S_iE_\lambda S_kE_\mu
   =S_iS_kE_\lambda E_\mu       \\
  &=S_\ell E_{\lambda'}E_\lambda E_\mu
   =S_\ell E_\alpha.
\end{split}
\]
Thus the transformations (33) form a group. The mixed group in question therefore depends on a continuous group $G_r$ and on $m$ fixed transformations $S_1,S_2,\ldots,S_m$ permutable with $G_r$ and such that the product of any two of them is equal to a third, multiplied by a transformation of the group $G_r$.

The elements in the formation of a mixed group depend on the groups $G_r$, having first the transformations $S_i$ permutable with $G_r$, a problem with which we shall be concerned later. Secondly, one must choose $m$ such transformations $S_1,S_2,\ldots,S_m$ in such a way that, with respect to the group $G_r$, they give rise to an ordinary finite group. If the number of the $\infty^r$ series is composite, the mixed group is only discrete but infinite; the result is wholly analogous, but then the multiplier transformations
\[
 S_1,
 \quad S_2,
 \quad S_3,
 \quad \ldots
\]
are infinite in number and constitute a discontinuous group. In every case we see, as was asserted in \S{}4, that these mixed groups partake at once of the nature of the continuous groups $G_r$ and of that of discontinuous groups.

Apparently the mixed groups that we have considered here are of a special nature because the series that compose them have the same number $r$ of parameters; but in reality we see easily that the more general case reduces to this. Indeed, suppose that the $m$ series composing the group have respectively
\[
 r_1,r_2,\ldots,r_m
\]
parameters and let $r$ be the maximum of the numbers $r_i$. Combining a series with $r_i$ parameters with another having $r_j$ parameters, one obtains a single series the number of whose essential parameters is at most $r_i+r_j$ and, supposing $r_i\ge r_j$, at least $r_i$. Therefore those among the $m$ series which have the maximum number $r$ of parameters already constitute by themselves a mixed group $\Gamma$ of the kind considered above, and in particular there is in $\Gamma$ a subgroup $G_r$ (one particular series)
\[
 \text{a)}\qquad x'_i=f_i(x_1,x_2,\ldots,x_n; a_1,a_2,\ldots,a_r)
\]
containing the identity. But then, if we consider the substitutions of another of the series with $r_i<r$ parameters,
\[
 \text{b)}\qquad x''_i=\varphi_i(x'_1,x'_2,\ldots,x'_n; b_1,b_2,\ldots,b_{r_i}),
\]
the composite series has $r$ parameters and therefore is in $\Gamma$. But in this series b) is completely contained, because in a) there is the identity. Thus one may simply suppress the series b) as already contained in another one of $\Gamma$.

From the preceding investigations it follows that the study of finite continuous groups reduces in substance solely to that of the groups $G_r$ generated by $r$ infinitesimal transformations
\[
 X_1f,
 \quad X_2f,
 \quad \ldots,
 \quad X_rf,
\]
and containing the identity and, with every transformation, its inverse. Henceforth, therefore, we shall concern ourselves only with these groups $G_r$, on which all the interest is concentrated.

\chapter*{Chapter III}
\addcontentsline{toc}{chapter}{Chapter III}

\begin{center}
\begin{minipage}{0.88\linewidth}
\centering
Equations of definition of a group $G_r$.\par
Order of the infinitesimal transformations.
\end{minipage}
\end{center}

\section*{\S{} 33. Elimination of arbitrary constants from a system of functions}

Every group $G_r$ is perfectly defined by its $r$ infinitesimal transformations
\[
 X_k f=\sum\limits_i^{1,\ldots,n}\xi_{ki}(x_1,x_2,\ldots,x_n)\frac{\partial f}{\partial x_i},
\]
satisfying the composition conditions
\[
 (X_kX_l)=\sum\limits_s^{1,\ldots,r}c_{kls}X_sf.
\]
The most general infinitesimal transformation of such a $G_r$ is given by
\[
 Xf=\sum\limits_k^{1,\ldots,r}e_kX_kf
   =\sum\limits_k^{1,\ldots,r}\sum\limits_i^{1,\ldots,n}
     e_k\xi_{ki}\frac{\partial f}{\partial x_i},
\]
where $e_1,e_2,\ldots,e_r$ are $r$ arbitrary constants. Since the coefficients $\xi_i$ of such a transformation $Xf$ are given by
\[
 \xi_i=\sum\limits_k^{1,\ldots,r}e_k\xi_{ki}
 \qquad (i=1,2,\ldots,n),
\]
they contain (linearly and homogeneously) the $r$ arbitrary constants $e_i$.

Once the $\xi_i$ are fixed, the group $G_r$ is determined, and hence so are the expressions of the $\xi_i$. But instead of giving the $\xi_i$ explicitly, one may also define them as the most general solutions of a certain system of simultaneous partial differential equations, which have great importance in many investigations and are called the equations of definition of the group.

In order to arrive at the construction of the system in question, we must first put forward some general observations on the elimination of constants from a group of functions of any number of independent variables.

Let $z_1,z_2,\ldots,z_m$ be $m$ functions of $n$ independent variables $x_1,x_2,\ldots,x_n$, also containing $r$ parameters (or constants) $a_1,a_2,\ldots,a_r$, which we may suppose, and shall suppose, essential, namely
\begin{equation*}\tag{1}
 z_i=f_i(x_1,x_2,\ldots,x_n; a_1,a_2,\ldots,a_r),
 \qquad i=1,2,\ldots,m.
\end{equation*}
Since in what follows we shall have to consider the derivatives of the various orders of the $z$, it is useful to introduce an abbreviated notation to denote them, and we shall put
\[
 \frac{\partial z_i}{\partial x_k}=z'_{i;k},
 \qquad
 \frac{\partial^2 z_i}{\partial x_k\partial x_h}=z'_{i;kh},
 \qquad \ldots,
\]
in general
\[
 \frac{\partial^s z_i}{\partial x_{k_1}\partial x_{k_2}\cdots\partial x_{k_s}}
   =z^{(s)}_{i;k_1k_2\ldots k_s}.
\]
This being set, we add to (1) the equations which result from it by a first differentiation,
\begin{equation*}\tag{$1_1$}
 z'_{i;k}=f'_{i;k}(x,a)
 \qquad
 \begin{gathered}
 i=1,2,\ldots,m,\\
 k=1,2,\ldots,n,
 \end{gathered}
\end{equation*}
then those obtained from it by a second differentiation,
\begin{equation*}\tag{$1_2$}
 z'_{i;k_1k_2}=f'_{i;k_1k_2}(x,a),
 \qquad \ldots,
\end{equation*}
and so we continue, prolonging the differentiation until the elimination of the $r$ parameters from the system of equations
\[
 (1),(1_1),(1_2),\ldots
\]
is possible, which must necessarily occur, as we shall now specify. Taking up again for this purpose the notation of \S{}3, let $\mu_0$ denote the characteristic of the functional matrix of (1), that is, of the $f_i$ with respect to the parameters; let $\mu_1$ denote that of the system $(1),(1_1)$, etc., namely of the functions $f_i,f_{i;k},f_{i;kh}$, and so on. From what we have seen in \S{}3, the numbers
\[
 \mu_0,\mu_1,\mu_2,\ldots,
\]
which are all $\le r$, form an increasing series up to the maximum $=r$ (because the $r$ parameters $a$ are, by hypothesis, essential); and from that point on the series of the $\mu$ is stationary. Let $\mu_{s-1}$ be the first $\mu=r$ and we shall have
\[
 \mu_0<\mu_1<\cdots<\mu_{s-2}<\mu_{s-1}=r.
\]

The system of equations
\[
 (1),(1_1),(1_2),\ldots,(1_{s-1})
\]
whose functional matrix of the right-hand sides with respect to the $r$ parameters $a$ has precisely characteristic $r$, is certainly solvable with respect to these parameters $a$, which will thus be expressed in terms of the $x$, the $z$, and their derivatives up to order $s-1$. These values of the $a$, substituted in the equations themselves $(1),(1_1),\ldots,(1_{s-1})$, give rise to a number $q$ of finite equations among the $x$, the $z$, and their derivatives up to order $s-1$:
\begin{equation*}\tag{$\alpha$}
 W_j(x,z,z'_{i;k},\ldots,z^{(s-1)}_{i;k_1\ldots k_{s-1}})=0,
 \qquad j=1,2,\ldots,q.
\end{equation*}
This number $q$ equals the difference between the total number of equations $(1),(1_1),\ldots,(1_{s-1})$ and the number $r$. If, in general, $\varepsilon_k$ denotes the number of derivatives of the $z$ of order $k$\begingroup\renewcommand\thefootnote{(*)}\footnotemark\endgroup{} and, for symmetry, we put $\varepsilon_0=m$, we shall have
\begingroup\renewcommand\thefootnote{(*)}\footnotetext{Evidently $\varepsilon_k=m\dfrac{n(n+1)\cdots(n+k-1)}{k!}$.}\endgroup{}
\begin{equation*}\tag{2}
 q=\varepsilon_0+\varepsilon_1+\cdots+\varepsilon_{s-1}-r.
\end{equation*}

Let us also consider the successive equations $(1_s)$, namely
\[
 z^{(s)}_{i;kk_1\ldots k_s}=f^{(s)}_{i;kk_1\ldots k_s}(x,a),
\]
and put in place of the $a$ their expressions in terms of the $x$, the $z$, and the derivatives up to order $s-1$; we shall thus have equations of the type
\begin{equation*}\tag{$\beta$}
 z^{(s)}_{i;kk_1\ldots k_s}
 =\psi^{(i,kk_1\ldots k_s)}(x,z,z'_{i;k},\ldots,z^{(s-1)}_{i;k_1\ldots k_{s-1}}),
\end{equation*}
by means of which all the partial derivatives of order $s$ of the $z$ are expressed in terms of those of lower order, of the $z$, and of the $x$. The system of partial differential equations $(\alpha),(\beta)$ is precisely the one that we proposed to construct.

\section*{\S{} 34. The system $(\alpha),(\beta)$ as a Mayer system. Properties of these systems}

We immediately recognize some important properties of this system, recalling that it is satisfied by the given functions
\[
 z_i=f_i(x_1,x_2,\ldots,x_n; a_1,\ldots,a_r)
\]
containing $r$ essential arbitrary constants. If, besides the $z_i$, we take as just as many new unknown functions the derivatives of the $z$ up to order $s-1$, we can substitute for the system $(\alpha),(\beta)$ a perfectly equivalent system of total differential equations for the $\varepsilon_0+\varepsilon_1+\cdots+\varepsilon_{s-1}$ unknown functions
\[
 z_i,
 \quad z'_{i;k},
 \quad z'_{i;k_1k_2},
 \quad \ldots,
 \quad z^{(s-1)}_{i;k_1\ldots k_{s-1}},
\]
combined with the $q$ finite equations among these that result from $(\alpha)$.

Since, by the theory of total differential equations, the general integral system can contain at most
\[
 \varepsilon_0+\cdots+\varepsilon_{s-1}-q=r
\]
arbitrary constants and will contain precisely $r$ only when the differential system is unlimitedly integrable. But this maximum $r$ is actually attained by the expressions $z_i=f_i(x,a)$ (and their derivatives up to order $s-1$), whence one concludes that these expressions give us the general integral of the differential system, which must be unlimitedly integrable.

After this we see that the system of partial differential equations $(\alpha),(\beta)$ enjoys the following properties:

a) The system permits all partial derivatives of order $s$ of the unknown functions $z$ to be expressed in terms of derivatives of lower order, of the $z$, and of the $x$ themselves. It is impossible, however, to express all derivatives of order $s-1$ in terms of those of lower order\begingroup\renewcommand\thefootnote{(*)}\footnotemark\endgroup.
\begingroup\renewcommand\thefootnote{(*)}\footnotetext{For this to be possible, indeed, it would be necessary that from the equations $(\alpha),(1),\ldots,(1_{s-2})$ the $r$ $a$ could be eliminated so as to substitute them in $(1_{s-1})$; but this is not so because $\mu_{s-2}<r$.}\endgroup{}

b) Differentiating $(\alpha),(\beta)$ and substituting in them, for the derivatives of order $s$, the expressions $(\beta)$, and eliminating those of order $s+1$, one obtains only equations that are consequences of $(\alpha),(\beta)$ themselves\begingroup\renewcommand\thefootnote{(**)}\footnotemark\endgroup. In another way, one may say that the differential consequences of the system $(\alpha),(\beta)$ up to order $s+1$ are already contained in the system itself.
\begingroup\renewcommand\thefootnote{(**)}\footnotetext{This follows from the unlimited integrability of the system of total differential equations.}\endgroup{}

In general, a system of simultaneous partial differential equations which enjoys the properties a), b), is called a Mayer system, because its integration is reduced to that of a system of unlimitedly integrable total differential equations, systems which Mayer especially considered. Indeed, if one eliminates as so many new unknown functions the first, second, etc. derivatives up to order $s-1$ of the $z$, one may substitute a system of total differential equations which, by property b), is unlimitedly integrable. The actual number of arbitrary constants which the general integral contains is always given by
\[
 r=\varepsilon_0+\varepsilon_1+\cdots+\varepsilon_{s-1}-q.
\]

From the developments of the preceding \S{} it follows in general that: If a system of simultaneous partial differential equations is such that only a finite number of arbitrary constants enters the expression of its general integral, it is equivalent to a Mayer system.

Often it is useful to substitute for the equations $(\alpha)$ of a Mayer system an equivalent system obtained as follows. It may happen that among the equations of the system $(\alpha)$ there are some containing only the $x,z$, or that there are such equations which are consequences of the system $(\alpha)$. Suppose that there exist precisely $\nu_0$ relations among only the $x,z$ as consequences of the system $(\alpha)$. To recognize the value of $\nu_0$ it suffices to compute the characteristic of a matrix, namely the functional matrix of $W_1,W_2,\ldots,W_q$ with respect to
\[
 z'_{i;k},
 \quad z'_{i;k_1k_2},
 \quad \ldots,
 \quad z^{(s-1)}_{i;k_1\ldots k_{s-1}}.
\]
If this characteristic is given by $\sigma_0$, which is necessarily $\le q$, then
\[
 \nu_0=q-\sigma_0.\begingroup\renewcommand\thefootnote{(*)}\footnotemark\endgroup
\]
\begingroup\renewcommand\thefootnote{(*)}\footnotetext{The fact that $\sigma_0$ is that characteristic means that among the equations $W_1,\ldots,W_q$, the $x$ and the $z$ there subsist $q-\sigma_0$ independent relations, which, for $W_1=0,\ldots,W_q=0$, give just as many relations among the $x$ and the $z$.}\endgroup{}
Similarly, let $\nu_1$ denote the number of relations among
\[
 x,z,z'_{i;k}
\]
which follow from $(\alpha)$, not counting the preceding $\nu_0$ among the $x,z$; let $\nu_2$ denote the number of relations consequent on $(\alpha)$ among the $x,z$ and the first and second derivatives. The numbers
\[
 \nu_0,\nu_1,\nu_2,\ldots,\nu_{s-1}
\]
will be calculated like the first $\nu_0$, by observing the values of the characteristics of functional matrices, and one will naturally have
\[
 \nu_0+\nu_1+\nu_2+\cdots+\nu_{s-1}=q,
\]
and hence
\[
 r=(\varepsilon_0-\nu_0)+(\varepsilon_1-\nu_1)+\cdots+(\varepsilon_{s-1}-\nu_{s-1}).
\]
For the system $(\beta)$ one can substitute an equivalent system formed, first, by $\nu_0$ equations expressing $\nu_0$ of the $z$ in terms of the remaining $z$ and the $x$; second, by $\nu_1$ equations giving $\nu_1$ of the first derivatives of the $z$ expressed in terms of the remaining $z'$, the $z$, and the $x$, etc.; in general $\nu_k$ of the derivatives of order $k$ will be expressed in terms of the remaining $\varepsilon_k-\nu_k$ and those of lower order $(k=0,1,2,\ldots,s-1)$. It is clear that one will always have
\[
 \nu_k<\varepsilon_k,
\]
for if all the $\varepsilon_k$ derivatives of order $k$ could be expressed in terms of those of lower order, we would have, for $k<s$, this would contradict the first a) of the properties of the Mayer system.

\section*{\S{} 35. The equations of definition of a group $G_r$}

We now apply the general observations of the two preceding \S\S{} to the formation of that system of partial differential equations (of Mayer type) which is satisfied by the coefficients $\xi_i$ of the most general infinitesimal transformation
\[
 Xf=\sum\limits_i^{1,\ldots,n}\xi_i\frac{\partial f}{\partial x_i}
\]
of the group $G_r$ generated by the $r$ infinitesimal transformations
\[
 X_k f=\sum\limits_i^{1,\ldots,n}\xi_{ki}(x_1,\ldots,x_n)\frac{\partial f}{\partial x_i}.
\]
One has
\begin{equation*}\tag{3}
 \xi_i=e_1\xi_{1i}+e_2\xi_{2i}+\cdots+e_r\xi_{ri},
 \qquad i=1,2,\ldots,n,
\end{equation*}
and the $\xi$ contain linearly and homogeneously the $r$ essential arbitrary constants $e_1,e_2,\ldots,e_r$. According to what was said in general in \S{}33, to (3) one associates the equations
\begin{equation*}\tag{$3_1$}
 \frac{\partial\xi_i}{\partial x_k}
  =e_1\frac{\partial\xi_{1i}}{\partial x_k}
  +e_2\frac{\partial\xi_{2i}}{\partial x_k}
  +\cdots
  +e_r\frac{\partial\xi_{ri}}{\partial x_k},
 \qquad i,k=1,2,\ldots,n.
\end{equation*}


\end{document}
